% qrh-main-remainder-height-v1
% qrh-direct-exposition-v2
% qrh-direct-exposition-v1
% qrh-natural-full-density-v1
% qrh-source-alignment-v1
\hypertarget{sources:chapter-25.-squarefree-preprocessing-and-ordinary-recovery}{%
\section{Chapter 25. Squarefree preprocessing and ordinary recovery}\label{sources:chapter-25.-squarefree-preprocessing-and-ordinary-recovery}}

\hypertarget{sources:final-route-coefficient-scope-and-uniform-choices}{%
\section{Final route, coefficient scope, and uniform choices}\label{sources:final-route-coefficient-scope-and-uniform-choices}}

This note fixes the route used by Chapters 25--30 of the complete proof. The accepted earlier theorem, weakened as stated above, supplies the fixed input \[
 \theta_I=0.87481891,\qquad \kappa_I=0.74963782,\qquad
 K_I=6\kappa_I=4.49782692.
\] The new endpoint is not used in proving any member of the source family.

\hypertarget{sources:the-actual-coefficient-class}{%
\subsection{The actual coefficient class}\label{sources:the-actual-coefficient-class}}

The ordinary source theorem concerns two factorized smooth ordinary ideal sums, with the literal characters, natural zeros, common punctures, fixed-inducing exclusions, finite profile seminorms, and height conventions of the inherited source identities. It includes all shorter plain lengths under its stated total affine cap.

The temporary ordinary restriction is \(\mu^2(n_1n_2)\). Retain also the squarefree projector on the complete marked product during common classification. The canonical marked prime pools may overlap; their exact collision decomposition is proved below. Both restrictions remain in the transformed coefficient. They are removed at completed positive norms by the divisor identities and the finite slot-cardinality reconstruction.

The allowed new operations are the two exact divisor kernels, actual extracted common punctures, outside character scalars of modulus at most one, and the particular finite Newton correction of the amplifier. Bounded pointwise multiplication of a coefficient is not substituted for a quadratic-norm argument.

\hypertarget{sources:elementary-centering-is-the-final-premise}{%
\subsection{Elementary centering is the final premise}\label{sources:elementary-centering-is-the-final-premise}}

The final route uses the contour \[
 \Re s_i=\frac{501}{1000},\qquad
 \alpha=\frac{499}{1000},\qquad
 L_{\mathrm{ref}}(b)=\frac{49}{100}-K_I b.
\] The centered terminal gate is \[
 \frac56+\frac{499}{1000}L_{\mathrm{ref}}(b)-T-9a>0.
\] The exact divisor transport below proves this saving; the global source theorem uses it in (G5). The finite source schedule uses this elementary saving and density \(\nu\ge1/2\).

\hypertarget{sources:one-mesh-works-through-the-whole-finite-induction}{%
\subsection{One mesh works through the whole finite induction}\label{sources:one-mesh-works-through-the-whole-finite-induction}}

Every application of an earlier source theorem occurs at an actual ambient scale \(U'=U^Q\), where the source certificate and common geometry give a uniform strictly positive lower bound for \(Q\). In particular, the proof uses \(Q\ge F-T/2-o(1)>0\).

For a path \(\gamma\) in the finite source graph, let \(q_\gamma>0\) be the product of uniform lower bounds for these scale exponents along the path. The empty path has \(q_\gamma=1\). At every source node choose positive local thresholds for the marked slot mesh, amplifier-pool separation, greedy-selection error, allowed height exponent, and logarithmic support erosion, inside the strict reserves of that node. Then choose the corresponding outer thresholds smaller than the minimum of \(q_\gamma\) times those local thresholds, over all the finitely many paths and nodes.

That minimum is positive. If a slot or height has exponent \(\tau\) in the original scale, its exponent at a child scale is \(\tau/Q_\gamma\); the preceding choice makes it smaller than the child threshold. The same calculation applies to support erosion. In particular every amplifier pool lies above each marked slot support at every recursive source invocation. The mesh choices precede even the ambient finite slot type. Once that type and the fixed profiles are chosen, the collision constants, finite controls and polynomial height degrees are fixed before moving arithmetic labels and heights. The eventual scale and large-height thresholds follow the required degrees.

The displayed \(10^{-12}\) mesh and padding values in the exact certificate are upper accounting allowances. Actual local and outer parameters are chosen finer after the finite positive reserves and transform constants are fixed. Thus the pullback just described requires no additional numerical inequality and does not assume that an unmeasured analytic loss equals \(10^{-12}\).

\hypertarget{sources:squarefree-total-product-preprocessing-for-a-two-plain-source}{%
\section{Squarefree total-product preprocessing for a two-plain source}\label{sources:squarefree-total-product-preprocessing-for-a-two-plain-source}}

\hypertarget{sources:status-and-purpose}{%
\subsection{Status and purpose}\label{sources:status-and-purpose}}

This note proves exact coefficient identities and weighted mass bounds. It also specifies the restricted coefficient class that must remain intact through the two source transforms, and proves its elementary extraction law. The source coefficient is the finite convolution (160), with its exact transforms (166)--(168). The present section proves the two divisor identities and the restriction retained by those transforms.

The central distinction is timing. The squarefree-product restriction is retained while the first common ideals are extracted. Its divisor expansion is used only after the complete transformed positive row norm has been reconstructed and a previously proved natural or initial moment is to be invoked.

Throughout, ideals are good integral ideals over the same fixed Eisenstein field. A presented multiplicative character retains all of its natural zeros. The fixed bad-prime set may be enlarged once, before scale, to contain all prime ideals of norm below 16. This is already an allowed operation in the supplied physical probe.

\hypertarget{sources:the-exact-two-variable-projector}{%
\subsection{1. The exact two-variable projector}\label{sources:the-exact-two-variable-projector}}

For nonnegative integers i,j define

\[
 f(i,j)=\mathbf 1_{i+j\le1}.
\]

Its two-variable generating series is 1+X+Y. The ordinary unconstrained pair has series 1/((1-X)(1-Y)). Define the multiplicative two-ideal coefficient h by

\[
 h_p(X,Y)=(1-X)(1-Y)(1+X+Y)
 =1-X^2-XY-Y^2+X^2Y+XY^2.                 \tag{1}
\]

Then, for all ideals a,b,

\[
 \mu^2(ab)=\sum_{d_1\mid a}\sum_{d_2\mid b}h(d_1,d_2).       \tag{2}
\]

Here mu-squared(ab) is one exactly when both a,b are squarefree and coprime. Equation (2) follows prime by prime: multiplication of (1) by 1/((1-X)(1-Y)) gives 1+X+Y. In particular it is a finite, exact divisor identity, with no analytic convergence issue.

The only nonzero local h values, apart from h\_\allowbreak{}p(0,0)=1, are

\begin{longtable}[]{@{}lr@{}}
\toprule\noalign{}
Exponents of d1,d2 & h\_\allowbreak{}p \\
\midrule\noalign{}
\endhead
\bottomrule\noalign{}
\endlastfoot
(2,0) & -1 \\
(1,1) & -1 \\
(0,2) & -1 \\
(2,1) & 1 \\
(1,2) & 1 \\
\end{longtable}

Let k be the convolution inverse of h in the two-ideal divisor algebra:

\[
 k_p(X,Y)=h_p(X,Y)^{-1}.                         \tag{3}
\]

Its coefficients are integers, k(1,1)=1, and it has no local monomials of total degree one. Thus

\[
 1=\sum_{d_1\mid a}\sum_{d_2\mid b}
 k(d_1,d_2)\mu^2((a/d_1)(b/d_2)).                \tag{4}
\]

Both (2) and (4) are coefficient identities for each finite pair (a,b).

\hypertarget{sources:subpower-normalized-mass-in-both-directions}{%
\subsection{2. Subpower normalized mass, in both directions}\label{sources:subpower-normalized-mass-in-both-directions}}

Fix sigma\textgreater1/2 and put z\_\allowbreak{}p=(Np)\^{}(-sigma). The absolute local h series is

\[
 1+3z_p^2+2z_p^3.
\]

The geometric expansion of (3) gives the absolute coefficient bound

\[
 \sum_{i,j\ge0}|k_p(i,j)|z_p^{i+j}
 \le\frac1{1-3z_p^2-2z_p^3}.                   \tag{5}
\]

For Np\textgreater=16, the denominator is bounded positively, already at sigma=1/2: 3/16+2/64=7/32\textless1. Both Euler products converge absolutely when sigma\textgreater1/2, because the degree-two prime sum has exponent 2sigma\textgreater1. This uses only the supplied absolute ideal Euler convergence to the right of one. The finitely many excluded primes remain natural zeros.

Consequently, for each epsilon\textgreater0 and each fixed B0,

\[
 \sum_{Nd_1Nd_2\le U^{B_0}}
 \frac{|h(d_1,d_2)|+|k(d_1,d_2)|}{\sqrt{Nd_1Nd_2}}
 \ll_{\epsilon,B_0} U^\epsilon.                \tag{6}
\]

Indeed choose sigma=1/2+epsilon/B0 when B0\textgreater0 and apply Rankin's inequality to the convergent products. The B0=0 case is bounded. Divisor weights of any fixed order are permitted by reducing the Rankin margin first; the resulting local series still has no degree-one term and its degree-two coefficients are bounded constants.

\hypertarget{sources:exact-transfers-of-centrally-normalized-products}{%
\subsection{3. Exact transfers of centrally normalized products}\label{sources:exact-transfers-of-centrally-normalized-products}}

For a natural presented character psi define

\[
 S_\psi(X;W)=X^{-1/2}\sum_a\psi(a)W(Na/X),
\]

and

\[
 T_\psi(X_1,X_2;W_1,W_2)
 =(X_1X_2)^{-1/2}\sum_{a,b}\mu^2(ab)\psi(ab)
 W_1(Na/X_1)W_2(Nb/X_2).
\]

Equations (2) and (4) imply, exactly,

\[
 T_\psi(X_1,X_2)
 =\sum_{d_1,d_2}\frac{h(d_1,d_2)\psi(d_1d_2)}
 {\sqrt{Nd_1Nd_2}}
 S_\psi(X_1/Nd_1)S_\psi(X_2/Nd_2),              \tag{7}
\]

\[
 S_\psi(X_1)S_\psi(X_2)
 =\sum_{d_1,d_2}\frac{k(d_1,d_2)\psi(d_1d_2)}
 {\sqrt{Nd_1Nd_2}}
 T_\psi(X_1/Nd_1,X_2/Nd_2).                     \tag{8}
\]

The fixed annular supports make the sums finite. Each scale is only shortened. If the character has a common puncture R, multiplicativity with zeros gives the identical formulas using psi\_\allowbreak{}R everywhere. No nonunit value is replaced by one.

An arbitrary fixed product of the original marked prime polynomials may multiply every term of (7) or (8). The factor psi(d1d2), though row dependent, has modulus at most one and is outside the child two-plain polynomial. Minkowski in the same positive row Hilbert space, followed by (6), therefore transfers a U\^{}(Q+epsilon) second-moment bound in either direction, provided that every shortened child is admitted. There is no arbitrary moving character in this operation: the child character is exactly psi, and the divisor character is a contracting row scalar.

Every divisor summand in (7) and (8) shortens each ordinary factor. The global domain (G1) admits all such shorter pairs under the same total affine cap. The common-mask transport at the moving puncture \(\mathfrak R d_1d_2\), proved below, pays its radical small-power cost uniformly before that puncture and the heights. Choose both this loss and the loss in (6) inside the finite strict reserves.

The same formulas hold for the centered difference of equal-product rectangles. Each divisor pair divides both rectangles by the same norms, so their product scales remain equal. At a fixed-inducing terminal one retains that difference until cancellation has been applied. The ordinary equal-product bound is the minimum of the lattice absolute bound \(O(\sqrt T)\) and the centered bound \(O(\sqrt T/L)\), including all positive subunit scales. Thus, for \(0\le\alpha\le1\), it is \(O(\sqrt T L^{-\alpha})\). After the same divisor extraction it costs \((Nd_1Nd_2)^{-(1-\alpha)}\). Take \(\alpha=499/1000\), for which \(1-\alpha=501/1000>1/2\), and apply the exact forward-kernel mass. The finite Newton correction is transported through the same complete centered coefficient, with its fixed finite mass. This proves the elementary terminal saving without treating bounded pointwise multiplication as contraction of a quadratic norm.

\hypertarget{sources:exact-extraction-closure-of-the-restriction}{%
\subsection{4. Exact extraction closure of the restriction}\label{sources:exact-extraction-closure-of-the-restriction}}

For fixed extracted ordinary factors A,B and variable quotients a,b,

\[
 \mu^2(ABab)
 =\mu^2(AB)\mu^2(ab)\mathbf1_{(ab,AB)=1}.       \tag{9}
\]

If AB is not squarefree, the whole allocated coefficient is zero. If AB is squarefree, the remaining coefficient is again the squarefree total-product projector, together with the additional common puncture rad(AB). This is a statement about the literal original ordinary factors, and it is independent of the choice of common slot allocation.

Repeated first and second common extractions replace A,B by their accumulated extracted factors. Equation (9) applies to the accumulated product and retains its cross-prime condition. It follows that the restricted ordinary coefficient class is closed under these exact divisor extractions. The ordinary quotients remain squarefree and mutually coprime until (7) is deliberately invoked at the final positive moment norm.

Row Poisson summation itself is linear in the column coefficients and does not change their mu-squared factor. The first complete common support extraction and the subsequent full artificial-coprimality Möbius identity therefore remain exact on this restricted class. Common Fourier separation touches only the smooth norm profiles, leaving (9) as an arithmetic restriction. No restriction may be discarded during that separation or before common multiplicities are classified.

In both transforms (166) and (137), retain this same restriction on the surviving ordinary indices and the complete common puncture. Apply a prior moment only after reconstructing its completed positive row norm, as in (140)--(141).

\hypertarget{sources:marked-pool-collision-reduction}{%
\subsection{5. Marked-pool collision reduction}\label{sources:marked-pool-collision-reduction}}

The canonical marked prime pools may overlap. Let \(M\) be the live marked
product and \(K\) its frozen part. During the common classification retain
both \(\mu^2(n_1n_2)\) and \(\mu^2(KM)\). On a tuple of prime ideals,
\(\mu^2(M)\) is the indicator that its entries are distinct. For the finite
edge set \(\mathcal A=\{(i,j):i\ne j\}\),
\[
 \mathbf1_{\mathrm{injective}}(\mathfrak p_\bullet)
 =\prod_{(i,j)\in\mathcal A}
       (1-\mathbf1_{\mathfrak p_i=\mathfrak p_j})
 =\sum_{E\subseteq\mathcal A}(-1)^{|E|}
       \prod_{(i,j)\in E}\mathbf1_{\mathfrak p_i=\mathfrak p_j}.
 \tag{10a}
\]
For each equality pattern, its connected components determine the frozen
blocks. Every singleton keeps its full original pool and its literal
row character. For a nonsingleton block, sum the bounded coefficients by
Cauchy with their normalized square masses. The annular prime norm bounds
make each such mass bounded by a constant depending on its fixed block
size and coefficient bound. There are finitely many edge patterns.
Consequently, if every completed full-singleton-subfamily norm is at most
\(E\), finite Cauchy bounds the completed projected energy by
\(C_{N,B}E\), with \(N\) the ambient finite slot count and \(B\) a fixed
coefficient bound. This constant is independent of scale, live labels and
heights. Strong induction on the surviving slot count applies the same
decomposition to the collision terms and recovers the unrestricted
marked source from the retained-projector source estimates for every
surviving subset. All source calls use full singleton pools and the
original coefficients.

On a surviving marked term, \((K,M)=1\) and
\(\mu^2(KM)=\mu^2(K)\mu^2(M)\). Together with the ordinary projector this
bounds every complete column valuation by two. Thus every nonempty
genuine first common pair has
\[
 (i,j)\in\{(1,1),(1,2),(2,1),(2,2)\}.\tag{10}
\]
Artificial M\"obius conductors keep their separate provenance. The
sixth-power amplifier terms retain their original error payment.

\hypertarget{sources:first-common-definitions-and-their-shared-payment}{%
\subsection{6. First common definitions and their shared payment}\label{sources:first-common-definitions-and-their-shared-payment}}

Let \(C,D_c\) be the complete first common ideals, with widths \(c=\log_U NC\), \(d=\log_U ND_c\). Let \(\pi\) be their full common radical width, \(R\) the active radical width from (164)--(165), and \(L_c,L_d\) the numerator widths from (172). The quantities \(R_{0,c},R_{0,d}\) are the widths canceled when the corresponding local numerator exponent is zero. Set \[
 B_c=(3c-5d-R)_+/6,\quad B_d=(3d-5c-R)_+/6,\qquad
 \Sigma=(c+d-2\pi-R-B_c-B_d)/2\ge0.
 \tag{RC1}
\] Nonnegativity is (MF3.56). The first-transform outside exponent \(O\) and the two side allowances satisfy, after unit row normalization, \[
 O+(S_c+S_d)/2=1-\Sigma.\tag{RC2}
\] This is the common-support payment (127)--(129), retaining the same first allocation weights. Refine by the displayed common/numerator widths before Cauchy; fixed dyadic refinements cost only \(U^\epsilon\). For nonnegative sector majorants, \(\sqrt{N_cN_d}\le\sum_{\alpha,\beta} \sqrt{N_{c,\alpha}N_{d,\beta}}\). Subtract \(\Sigma\) once from each side before pairing. Fix an earlier face \((B,\Lambda)\). 

\subsection{7. Exact raw radius on this restricted class}

Let A,b,c,D now denote the nonnegative logarithmic widths of first common primes of types (1,1),(1,2),(2,1),(2,2), respectively. To avoid confusion, write C,Dc for the two total first common lengths in this paragraph. The exact local rules give

\[
 \begin{aligned}
 C&=A+b+2c+2D,&D_c&=A+2b+c+2D,\\
 R&=b+c,&L_C&=2c,&L_{D_c}&=2b,\\
 R_{0,C}&=b,&R_{0,D_c}&=c,&\Sigma&=D.
 \end{aligned}                                                    \tag{12}
\]

Both B positive parts vanish: 3C-5Dc-R and its reverse are sums of nonpositive local contributions. Hence

\[
 \tau_C=L_C-R-(C-D_c)_+=-(D_c-C)_+,
 \quad Q_C=F-(D_c-C)_+,                                          \tag{13}
\]

and similarly with C,Dc exchanged. In particular

\[
 Q_C\le F,\qquad
 |C-D_c|\le(C+D_c)/3.                                            \tag{14}
\]

Thus C+Dc\textless=u implies F-u/3\textless=Q\textless=F.

The individual RAW coefficient simplifies to

\[
 W_C=-b-D-l\{C-B(b-c)_+\}.                                      \tag{15}
\]

For 0\textless B\textless=1, C\textgreater=b implies W\_\allowbreak{}C\textless=0, individually. This statement does not eliminate a need to pay a HIGH partner's remaining coefficient; the paired HIGH and exceptional schedules still have to be checked.

The nonpositive bound (15) is used without any further frozen-slot deduction. The complete erasure cost, including both floors, is proved in (G10)--(G13).

\hypertarget{sources:chapter-26.-exact-corrected-amplification}{%
\section{Chapter 26. Exact corrected amplification}\label{sources:chapter-26.-exact-corrected-amplification}}

\hypertarget{sources:exact-amplification-with-a-corrected-squarefree-coefficient}{%
\section{Exact amplification with a corrected squarefree coefficient}\label{sources:exact-amplification-with-a-corrected-squarefree-coefficient}}

\hypertarget{sources:status-and-scope}{%
\subsection{Status and scope}\label{sources:status-and-scope}}

This is a new exact amplifier lemma for the restricted source class in the squarefree preprocessing section above. It uses the local Gauss identity (8), the first transform (120), and the second transform (137). It gives a single MAIN norm. The affine source face follows after the source inequalities and admission gates below are established.

The required condition is squarefreeness at amplifier primes. A column whose two ordinary factors have squarefree coprime product may still have valuation two at another prime contributed by a retained marked slot. The amplifier pool avoids all marked slot supports.

Fix a positive amplifier width \texttt{a}, the literal pool

\[
 \mathcal P=\{\mathfrak p:U^a/2<N\mathfrak p\le U^a,
                  \ \mathfrak p\text{ in the fixed permitted ray}\},
 \qquad P=\#\mathcal P.
\tag{CA1}
\]

The finite bad set fixed in Section 1.2 is omitted. Fix the pool before the centered rectangles, live labels and Mellin heights. The inherited fixed-extension prime estimate gives \texttt{P=U\^{}\{a+o(1)\}}. All norms and supports in the source are bounded by a fixed power of \texttt{U}. Consequently there is a fixed integer \texttt{R}, depending on that power and \texttt{a}, for which every supported column satisfies

\[
 r(u):=\#\{\mathfrak p\in\mathcal P:\mathfrak p\mid u\}\le R.
\tag{CA2}
\]

Increase the eventual threshold so that \texttt{P\textgreater{}2R+2}.

\hypertarget{sources:the-identity-at-every-frequency}{%
\subsection{1. The identity at every frequency}\label{sources:the-identity-at-every-frequency}}

Suppose \texttt{v\_\allowbreak{}p(u)} is zero or one for every pool prime. If \texttt{p} does not divide \texttt{u}, the primitive Gauss change of variable at \texttt{p} gives

\[
 \mathfrak g(u,p^6h)=\mathfrak g(u,h).
\]

If \texttt{p} divides \texttt{u} once, the local primitive Gauss sum at the frequency \texttt{p\^{}6h} is zero. If \texttt{p} also divides \texttt{h}, the local factor at \texttt{h} is already zero. Thus, including all nonunit and zero frequencies,

\[
 \mathfrak g(u,p^6h)
 =\mathbf1_{p\nmid u}\mathfrak g(u,h).
\tag{CA3}
\]

Let \texttt{b(u)} be the complete residual coefficient, including the full centered difference, common puncture, artificial divisibility projector, whole ray character and norm twist. Define

\[
 \widetilde b(u)=\frac{P}{P-r(u)}b(u),\qquad
 F_b(h)=\sum_u b(u)\mathfrak g(u,h).
\tag{CA4}
\]

Summing (CA3) over the entire pool yields the exact identity

\[
 \boxed{\quad F_b(h)=\frac1P
                  \sum_{p\in\mathcal P}F_{\widetilde b}(p^6h).
            \quad}
\tag{CA5}
\]

The coefficient on the right is independent of the individual prime \texttt{p}.

For a nonnegative annular majorant \texttt{Omega}, Cauchy's inequality gives

\[
 \sum_h\Omega(Nh/L)|F_b(h)|^2
 \le\frac1P\sum_{p,h}\Omega(Nh/L)
                      |F_{\widetilde b}(p^6h)|^2.
\tag{CA6}
\]

The exact pushforward through \texttt{(h,p)\ -\textgreater{}\ p\^{}6h} in the sixth-power amplification identity (131) gives: each nonzero output of polynomial norm has at most a fixed number of such prime preimages, because their product divides its norm. The annular majorant in (CA6) is supported away from zero. Any separate zero-frequency terminal is kept with its original support and terminal bound; no sparse pushforward count is asserted at zero. The smooth MAIN radius therefore majorizes the right side, with one density factor \texttt{P\^{}\{-1\}=U\^{}\{-a+o(1)\}} and a constant depending only on the fixed cells. Multiply by the original nonnegative allocation weights and sum; those weights are unchanged. The enlarged MAIN norm is transformed a second time with the whole coefficient \texttt{tilde\ b} still intact.

Use the bound \texttt{\textbar{}tilde\ b/b\textbar{}\textless{}=2} for diagonal and absolute terminal estimates. For the other completed positive norms use the exact finite divisor representation below.

\hypertarget{sources:exact-finite-divisor-representation}{%
\subsection{2. Exact finite divisor representation}\label{sources:exact-finite-divisor-representation}}

For integers \texttt{0\textless{}=r\textless{}=R}, Newton interpolation gives

\[
 \frac{P}{P-r}
 =\sum_{j=0}^{R}c_j(P)\binom rj,
 \qquad c_0(P)=1,\qquad
 c_j(P)=\frac{j!}{(P-1)(P-2)\cdots(P-j)}.
\tag{CA7}
\]

Terms with \texttt{j\textgreater{}r} vanish. The formula follows from the finite difference identity

\[
 \Delta^j\frac{P}{P-r}\bigg|_{r=0}
 =\frac{j!}{(P-1)\cdots(P-j)},
\]

and finite Newton interpolation.

After a complete second common extraction, let \texttt{r0} be the number of pool primes in the frozen extracted part. The remaining column is punctured at that part, so its pool count adds disjointly. The exact shifted version is

\[
 \frac{P}{P-r_0-r}
 =\sum_{j=0}^{R}c_j^{(r_0)}(P)\binom rj,
 \qquad
 c_j^{(r_0)}(P)=
 \frac{Pj!}{(P-r_0)(P-r_0-1)\cdots(P-r_0-j)}.
\tag{CA8}
\]

Here \texttt{R} can be enlarged once to cover all accumulated extractions. On the supported range \texttt{r0+r\textless{}=R}, all denominators are positive. For use in Mellin inversion, the finite polynomial on the right of (CA8) may be extended to all nonnegative \texttt{r}: the smooth supports ensure that only the range where it agrees with the desired multiplier contributes to the original sum.

For a column squarefree on the pool,

\[
 \binom{r(u)}j
 =\sum_{\substack{d\mid u,\ d\text{ squarefree}\\
                  p\mid d\Rightarrow p\in\mathcal P,\ \omega(d)=j}}1.
\tag{CA9}
\]

The corresponding divisors are kept as coefficient restrictions through first common classification and both row transforms. Equation (CA9) does not justify enlarging the first common multiplicity class.


\paragraph{Support of the actual Newton correction.}
Write \(P=\#\mathcal P\), the actual integer pool cardinality, and let
\(\Psi_z\) be the literal multiplicative row weight with accumulated
puncture \(\mathfrak R\). Keep the selected marked product \(M\) inside
the correction:
\[
 \frac{P}{P-r_{\mathcal P}(DMIJ)}\,
 \mu^2((A_1I)(A_2J))\Psi_z(IJ)
 \bigl(\text{the same centered profile}\bigr).
\]
Terms with \(\Psi_z(M)=0\) or \(\Psi_z(IJ)=0\) contribute zero. On
each surviving term, \(D\mid\mathfrak R\) gives \((D,MIJ)=1\).
Since the amplifier pool avoids the marked support,
\[
 r_{\mathcal P}(DMIJ)=r_{\mathcal P}(D)+r_{\mathcal P}(IJ).
\]
The support norm bound and prime lower norm bound give a fixed residual
depth \(J_0\). The actual pool budget
\[
 2\bigl(r_{\mathcal P}(D)+J_0+1\bigr)\le P
\]
then makes every supported denominator positive and the correction at
most two. Apply the finite shifted Newton expansion on these surviving
terms, retaining the same puncture, marked product, smooth profiles and
independent heights in both rectangles. Natural character zeros remain
in each summand.

\hypertarget{sources:restoration-at-completed-nonexceptional-row-norms}{%
\subsection{3. Restoration at completed nonexceptional row norms}\label{sources:restoration-at-completed-nonexceptional-row-norms}}

When a previously proved natural or initial affine moment is to be applied to a complete positive row norm, expand (CA8) and allocate each prime of \texttt{d} among the two ordinary factors. The pool avoids slots, so no marked prime has to be selected in this operation. There are at most \texttt{2\^{}j} allocations. Exact extraction divides both rectangles' first and second scales by the same assigned ideals, respectively, and contributes the central scalar

\[
 \frac{\psi(d)}{\sqrt{Nd}},\qquad |\psi(d)|\le1.
\tag{CA10}
\]

The remaining two ordinary factors still have squarefree coprime product and the accumulated common puncture. Thus their restriction can subsequently be restored using the finite divisor kernel \texttt{h} from the squarefree transfer lemma. All character zeros are retained. The whole row character is the same presented character; (CA10) is a contracting row scalar, not a replacement character.

For \texttt{P\textgreater{}=2(R+r0+1)}, the coefficients obey

\[
 c_j^{(r_0)}(P)\le 2^{j+1}j!P^{-j}.
\]

There are at most \texttt{P\^{}j/j!} squarefree choices of \texttt{d}, each with \texttt{Nd\ \textgreater{}=\ (U\^{}a/2)\^{}j}. Hence the total central absolute coefficient mass, including allocations, is bounded by

\[
 \sum_{j=0}^{R}c_j^{(r_0)}(P)
 \sum_{\substack{d\text{ as in (CA9)}\\\omega(d)=j}}
             \frac{2^j}{\sqrt{Nd}}
 \le 2\sum_{j=0}^{R}(4\sqrt2\,U^{-a/2})^j=O_R(1).
\tag{CA11}
\]

Minkowski in the same positive row Hilbert space now applies. This is a uniform finite mass calculation, with no new prime-label count after the completed norm. The separate \texttt{h} restoration costs only \texttt{U\^{}epsilon} as proved in the transfer lemma. Shortened plain factors must satisfy the relevant earlier affine admission or its proved natural fallback; the coefficient identity does not erase those gates.

% qrh-joint-newton-family-v1
\paragraph{One Newton bound for the complete row family.}
Use the same pool \(\mathcal P\), with \(P=\#\mathcal P>0\),
for the finite row family. The accumulated divisor \(D_u\), the
ordinary factors, puncture, character and surviving slot subset may
depend on the row \(u\). Choose the common depth \(J_0\) from their
polynomial norm caps. For \(S\subseteq\mathcal P\), \(j=|S|\le J_0\),
and \(A\subseteq S\), let \(\mathfrak d_{S,A}\) be the ordinary
divisor pair assigned by the exact Newton expansion, writing
\(N\mathfrak d_{S,A}=N\mathfrak d_1\,N\mathfrak d_2\). Its scalar
coefficient, including the natural character values, has absolute value
at most
\[
 w_{S,A}=\frac{2^{j+1}j!}{P^j\sqrt{N\mathfrak d_{S,A}}}.
\]
The pool budget above proves this for every \(D_u\) simultaneously.
Since \(N\mathfrak d_{S,A}\ge1\),
\[
 \sum_{\substack{S\subseteq\mathcal P\\|S|\le J_0}}
 \sum_{A\subseteq S}w_{S,A}
 \le M_{J_0}:=\sum_{j=0}^{J_0}2^{2j+1}j!.
\]
Indeed, there are \(\binom Pj\le P^j\) subsets of size \(j\)
and \(2^j\) assignments per subset.

Write the literal normalized corrected row as
\(H_u=\sum_{S,A}c_{S,A}(u)F_{S,A}(u)\), keeping its original
marked factor and both nonzero-row gates inside \(F_{S,A}\).
For any nonnegative original row weights \(\omega_u\), weighted
Cauchy followed by interchange of these finite sums gives
\[
 \sum_u\omega_u|H_u|^2
 \le M_{J_0}\sum_{S,A}w_{S,A}
          \sum_u\omega_u|F_{S,A}(u)|^2
 \le M_{J_0}^2 E
\]
whenever every displayed complete cofactor-family sum is at most \(E\).
Apply the ordinary singleton-source and HIGH count estimates to these
same simultaneous cofactor families, with their actual moving punctures
and common Mellin height. The finite factor \(M_{J_0}^2\) is independent
of the scale and moving row data.

% qrh-full-normalized-ratio-v1
\paragraph{Ratio restrictions of the full normalized family.}
Retain the complete nonnegative effective weight in each normalized
quadratic energy. Exact reindexing by the original row label shows
that restriction to any subtype or ratio fiber decreases this energy.
Let \(E_{\rm full}\) bound the full energy in both orientations,
with the stated polynomial Mellin-height envelope. Use this same
bound on every ratio fiber. The completed ratio contraction then gives
\[
 \left\|\sum_q\operatorname{phase}_q\operatorname{block}_q\right\|
 \le C\,\#\mathcal B_Z\,E_{\rm full},
\]
where \(\mathcal B_Z\) is the original finite set of ratio bins.
The constant is chosen before the moving row and label families.
For the actual complete family, the polynomial caps on its radial
scale, frequency and source radius imply
\(\#\mathcal B_Z\le C_\epsilon Z^\epsilon\).
Thus this contraction costs one arbitrarily small scale power and
retains both rays, residual projectors, phases and all original
outside weights. It introduces no source-label cardinality factor.

% qrh-first-allocation-family-v1
\paragraph{Selection of the actual source family.}
For an original first column \(j\), let \(\mathcal A_j\) be its
complete first-allocation set and write its complete nonnegative cost as
\[
 T_j=\sum_{B\in\mathcal A_j}|a_{j,B}|^2K_{j,B}.
\]
The cost \(K_{j,B}\) includes all subsequent blocks, residual-ideal
weights, windows, profiles and physical scalars. If \(T_j>0\), choose
an actual \(B_j\) maximizing the entire product
\(|a_{j,B}|^2K_{j,B}\). Then \(a_{j,B_j}\ne0\), \(K_{j,B_j}>0\), and
\[
 T_j\le\#\mathcal A_j\,|a_{j,B_j}|^2K_{j,B_j}.
\]
The zero fibers contribute zero. Retain the previous allocation-Cauchy
factor in the outside weight. Bound each of the two resulting allocation
cardinality factors by its ideal-divisor small-power estimate.

The full first-column code consists of its source label, seed, four
dyads and two ray characters. For a fixed source-label/seed pair, its
fiber therefore has cardinality at most
\(\#\mathcal D_Z\,\#\widehat{\mathcal R}^{\,2}\), where
\(\mathcal D_Z\) is the actual dyad envelope and
\(\widehat{\mathcal R}\) is the fixed finite ray-character group.
The polynomial bounds on the source radius, volume and reciprocal
normalizer give \(\#\mathcal D_Z\le C_\epsilon Z^\epsilon\).
Choosing a column of maximal complete cost in each nonzero fiber thus
leaves a family with injective source-label/seed key, at this uniform
small-power cost. Keep every actual corrected second block at its
chosen first allocation. Its modified coefficient, child character,
first seed, mask and original singleton slots are the literal ones of
the transformed source. All subsequent joint estimates use this family.

\paragraph{Forced labels and their actual allocations.}
For a complete allocation \(B\) of a common ideal \(C\), extraction sends
\[
 A_i\longmapsto A_iB_i,\qquad D\longmapsto DC,\qquad
 \mathfrak R\longmapsto\mathfrak RC.
\]
% qrh-joint-forced-window-v1
% qrh-fixed-residual-family-v1
The forced divisor seed remains present, and the invariants
\(A_1A_2\mid\mathfrak R\), \(D\mid\mathfrak R\) propagate.
Fix a squarefree nonzero forcing ideal \(L\). In the original finite
row set \(\mathcal U\), keep the rows for which \(L\) belongs to the
actual residual pool, using a membership indicator in a common finite
union of those pools. For each nonzero normalized mother row
\(H_{u,L}\), choose an actual complete allocation with nonzero frozen
coefficient and positive seed-one child norm. Its complete forced ideal
\(C_{u,L}\) satisfies \(L\mid C_{u,L}\mid L^2\), and its exact scalar is
\(\lvert a_B\rvert^2/\operatorname{rawReduction}(B)\).
Finite Cauchy and the ideal-divisor small-power bound give one
\(K_{N,\delta}>0\), before \(L\), the row family and all its moving
data, such that
\[
 |H_{u,L}|^2\le K_{N,\delta}(NL)^{\delta-1}
             \mathcal C_{u,L}|G_{u,L}|^2.
\]
Here \(\mathcal C_{u,L}\) is the full profile-dilation cost and
\(G_{u,L}\) retains the actual child volume, puncture and heights.
For nonnegative original weights \(\omega_u\), summation gives
\[
 \sum_{u\in\mathcal U}\omega_u|H_{u,L}|^2
 \le K_{N,\delta}(NL)^{\delta-1}
 \sum_{u\in\mathcal U_{L,+}}\omega_u
                   \mathcal C_{u,L}|G_{u,L}|^2.
\]
The live set \(\mathcal U_{L,+}\) is a subset of the original rows.
Fix the missing ideal, ray labels and common divisor pair before
applying the earlier source to this family. Since
\(L\mid C_{u,L}\mid L^2\), the prime supports of \(C_{u,L}\) and
\(L\) agree. The excluded-character mask therefore has the same
coefficients as the common mask with \(C_{u,L}\) replaced by \(L\).
This is one fixed puncture for the joint source estimate. Its constants
are uniform in \(L\). Partition this same fixed-\(L\) family into LOW
and HIGH, preserving its original characters, correction factors,
heights, masks, weights and first outside allocation scalar. After the
joint estimates, sum over the full residual union with its retained
inverse-norm weight.

For the two ordinary allocated ideals put \(d=N(B_1B_2)\ge1\). Every
live child slot scale is at least one, and the literal child data satisfy
\[
 V_{\rm child}\le V/d,\qquad
 L_{\rm child}\ge L_0/d,\qquad
 \mathcal S_{\rm child}\le\mathcal S,
\]
where \(L_0>0\) bounds the four original raw factors below and
\(\mathcal S\) is the original slot-control factor. For
\(0\le\alpha<1/2\),
\[
 \mathcal S_{\rm child}V_{\rm child}
       \max(1,L_0/d)^{-2\alpha}
 \le \mathcal S V\max(1,L_0)^{-2\alpha}.
\]
Indeed \(\max(1,L_0)\le d\max(1,L_0/d)\), and
\(d^{2\alpha-1}\le1\). All slot controls are nonnegative; multiplying
once more by \(\mathcal S_{\rm child}\le\mathcal S\) gives the same
inequality with squared slot controls when required by the energy bound.
% qrh-sharp-exceptional-allocation-v1
% qrh-selected-volume-geometry-v1
\paragraph{The full volume through the three extractions.}
Let \(b\ge1\) bound the fixed annular supports and let the original
marked set have at most \(N\) elements. For a contribution with
nonzero complete frozen coefficient, denote the two common ideals and
the complete forced ideal by \(\mathfrak C_1,\mathfrak C_2,
\mathfrak C_f\), and their actual allocations by \(B^{(1)},B^{(2)},
B^{(3)}\). The exact raw-scale identities give
\[
 V_3=\frac{X_1X_2\prod_{i\in T}P_i}
 {\prod_{k=1}^3\operatorname{rawReduction}(B^{(k)})}
 \le \frac{b^{3N}V_0}
 {N\mathfrak C_1\,N\mathfrak C_2\,N\mathfrak C_f},
 \qquad V_0=X_1X_2\prod_{i\in T}P_i.
\]
Every ordinary corner contracts separately, and the surviving singleton
product is the exact image subset of the original marked set.
For \(Z>1\), the logarithmic volume allowance is therefore
\(3N\log b/\log Z\). It is absorbed into the chosen fixed loss by
an eventual threshold depending on \(N\). The scalar normalization
and the four separate corner budgets retain their original factors.

% qrh-selected-slot-geometry-v1
\paragraph{The surviving original slot widths.}
The canonical marked profiles have lower support endpoint one. Keep
the three nonzero frozen coefficients in the preceding extraction.
Let \(T_0\) be the original marked set, \(F\subseteq T_0\) the
exact image of the final live slots, and \(P_i=Z^{w_i}\), with
\(w_i\ge0\). For a cutoff \(\omega_*\ge0\), write
\(F_+=\{i\in F:w_i\ge\omega_*\}\) and
\(\ell=\sum_{i\in F_+}w_i\). The frozen support bounds give
\[
 \sum_{i\in T_0}w_i
 -\log_Z(N\mathfrak C_1)-\log_Z(N\mathfrak C_2)
 -\log_Z(N\mathfrak C_f)-N\omega_*\le\ell.
\]
Thus, if \(\beta-\delta\le\sum_{i\in T_0}w_i\le\beta\) and
\(\delta+N\omega_*\le2\varepsilon_w\), then
\[
 0\le\ell\le\beta,\qquad
 \beta-2\varepsilon_w\le
 \ell+\log_Z(N\mathfrak C_1)+\log_Z(N\mathfrak C_2)
       +\log_Z(N\mathfrak C_f).
\]
For every collision singleton subset \(S\subseteq F\), put
\(s_+=\sum_{i\in S\cap F_+}w_i\). Nonnegativity gives
\[
 0\le s_+\le\ell,\qquad
 \sum_{i\in S}w_i+(\ell-s_+)\le\sum_{i\in F}w_i.
\]
These are the slot coordinates used with the full-volume bound above.

% qrh-masked-radical-allocation-v1
\paragraph{The common-ideal allocation payment.}
Let \(\mathfrak C\ne0\) be an actual common ideal,
\(\mathcal A_{\mathfrak C}\) its complete allocation set,
\(a_B\) the frozen coefficient and \(V_B\) the child volume.
Put \(d_B=N(B_1B_2)\) and
\(R_{\mathfrak C}=N\operatorname{rad}\mathfrak C\).
The ordinary squarefree projector makes the actual child row zero
unless \(B_1B_2\) is squarefree: every restored ordinary product
contains this frozen divisor. Retain this zero term in the original
allocation sum. On the surviving set \(\mathcal A_{\mathfrak C}^{\rm sf}\),
the squarefree divisor \(B_1B_2\mid\mathfrak C\) divides
\(\operatorname{rad}\mathfrak C\), so \(1\le d_B\le R_{\mathfrak C}\).
The exact raw-scale identity and frozen support give
\[
 |a_B|^2V_B^2\le\frac{\mathcal F^2}{(N\mathfrak C)^2}V^2,
 \qquad \max(1,L_0)\le d_B\max(1,L_0/d_B).
\]
Thus, for \(0\le\alpha<1/2\),
\[
\begin{split}
 |a_B|^2V_B^2
 \bigl[\max(1,L_0/d_B)^{-\alpha}\mathcal S_B\bigr]^2
 \le{}&\frac{\mathcal F^2d_B^{2\alpha}}{(N\mathfrak C)^2}
 V^2\bigl[\max(1,L_0)^{-\alpha}\mathcal S\bigr]^2.
\end{split}
\]
Cauchy with the complete original allocation cardinality gives
\[
 \#\mathcal A_{\mathfrak C}
 \sum_{B\in\mathcal A_{\mathfrak C}^{\rm sf}}
 \frac{\mathcal F^2d_B^{2\alpha}}{(N\mathfrak C)^2}
 \le (\#\mathcal A_{\mathfrak C})^2\mathcal F^2
       \frac{R_{\mathfrak C}^{2\alpha}}{(N\mathfrak C)^2}.
\]
The original frozen coefficients, profile-dilation cost, forcing seed,
first outside weight and row selector remain in the corresponding
completed norm.

% qrh-fixed-q-exceptional-v1
\paragraph{The fixed-modulus forcing count.}
Keep the fixed nonzero inducing modulus \(Q\) throughout the exceptional
row count. Split the squarefree forcing ideal as
\(F=F_{\rm good}F_Q\), according as its prime factors are coprime to
\(Q\) or divide \(Q\). The second factor divides
\(\operatorname{rad}Q\). At each prime of \(F_{\rm good}\), fixed
induction forces the row valuation into the residue class four modulo
six. The sixth-power parameterization therefore gives, on
\(N(z)\le C_rZ^M\),
\[
 \#\{z\text{ in the exceptional row set}\}
 \le C_Q C_r^{1/6}Z^{M/6}(NF_{\rm good})^{-2/3}
 \le C_Q(N\operatorname{rad}Q)^{2/3}
       C_r^{1/6}Z^{M/6}(NF)^{-2/3}.
\]
Here \(C_Q=768\,6^{\omega(Q)}\), with the fixed ray and unit
sectors retained. The second inequality follows from
\(NF\le NF_{\rm good}\,N\operatorname{rad}Q\). The modulus is fixed
before the source constant, so its radical factor is included in that
constant. The profile orders, source exponents and both independent
Mellin heights retain their stated quantifier order.

\paragraph{Paired windows and the full forcing sum.}
Let \(\mathfrak C,\mathfrak D\) be the two actual common ideals with
the same prime support, and put
\(R=N\operatorname{rad}\mathfrak C=N\operatorname{rad}\mathfrak D\).
Taking square roots of the two masked allocation sums gives
\[
 \mathcal W_{\mathfrak C}\mathcal W_{\mathfrak D}
 \le\#\mathcal A_{\mathfrak C}\#\mathcal A_{\mathfrak D}
       \frac{R^{2\alpha}}{N\mathfrak C\,N\mathfrak D}
       \mathcal F^2E_L\mathcal H_{J,t}^2.
\]
Here \(\mathcal H_{J,t}\) is the retained height envelope times its
Mellin profile moment. Each window uses its own residual puncture
and the same parent budget. For at most \(N\) marked slots, the
divisor bound uses all \(N+2\) slots and gives, for each fixed
\(\eta>0\),
\[
 \#\mathcal A_{\mathfrak C}\#\mathcal A_{\mathfrak D}
 \le A_{N,\eta}(N\mathfrak C\,N\mathfrak D)^\eta.
\]
% qrh-second-radical-scalar-v1
\paragraph{The second common radical payment.}
Write \(\mathcal R\) for the actual row radius, \(F_\nu\) for the
common Fourier-frequency ideal and \(\mathcal N\) for the literal
common normalizer. The radial majorant uses
\(H=\max(1,4\mathcal R/NF_\nu)\).
A nonzero physical block supplies a nonzero retained row, so
\(H=4\mathcal R/NF_\nu\) on that block. A zero block contributes zero.
At a common prime of norm \(q\), put
\(c=v_{\mathfrak p}(\mathfrak C)\),
\(d=v_{\mathfrak p}(\mathfrak D)\), and \(m=\min(c,d)\).
The canonical partition gives exponents \(u,v\) equal to one in
its equal-valuation unit and nonunit cases, respectively, and zero
otherwise. The forcing exponent is \(f=1\) in the \((1,1)\) nonunit
case and zero otherwise. Thus the prime exponents in
\(\mathcal N,NF_\nu,NF\) are \(m-u,m+v,f\).
The local factor of
\[
 \frac{\mathcal N}
 {N\mathfrak C\,N\mathfrak D\,(NF_\nu)^{1/6}(NF)^{2/3}}
\]
is \(q^{-\kappa}\), where
\[
 \kappa=c+d-\frac56m+u+\frac16v+\frac23f\ge2.
\]
For the four possible types, its lower bounds are respectively
\(2,13/6,13/6,5/2\). Multiplying over the common primes gives
\[
 \frac{\mathcal N H^{1/6}}
 {N\mathfrak C\,N\mathfrak D\,(NF)^{2/3}}
 \le (4\mathcal R)^{1/6}R^{-2}.
\]
Together with the paired window, the radical factor is
\(R^{2\alpha-2}\). The retained projectors give
\(N\mathfrak C\,N\mathfrak D\le R^4\), so the allocation-cardinality
loss costs at most \(R^{4\eta}\). For
\(\alpha=499/1000\) and \(\eta=1/10000\),
\(2\alpha+4\eta-2=-626/625<-1\).
The remaining divisor and finite type multiplicities are bounded by
an arbitrarily small power of \(R\); choose that power within this
strict surplus. The resulting inverse-radical weight is summed with
the ideal harmonic bound, with its arbitrarily small scale loss.
All original outside, dyadic and profile factors stay in the sum.

The literal corrected seed budget has the exact form
\(E_L=(NL)^{\varepsilon_L-1}E_1\).
% qrh-sharp-residual-pool-cap-v1
Use the norm cap \(NL\le Z^{B_L}\) for the full forcing pool.
For the actual second residual pool take \(B_L=4\), independently
of the number \(N\) of marked slots. Indeed, the four initial raw
scales are at most \(Z\). The product of the two stored ordinary
factor norms is at most a fixed support constant times \(Z^2\).
The complete marked product has norm at most
\(2^N Z^{\sum_i w_i}\le2^N Z\).
Thus every stored initial column product has norm at most
\(C_NZ^3\), and hence at most \(Z^4\) after an eventual threshold
depending on \(N\). Each quotient-child column reconstructs a stored
parent column after multiplication by its actual common allocation.
Its norm therefore contracts through both extractions. The actual
sector reconstruction gives \(L\mid J\), with \(\mathfrak D J\)
in that stored pool and \(\mathfrak D\ne0\). Hence
\(NL\le NJ\le N(\mathfrak D J)\le Z^4\) throughout the padded
residual union, including zero-coefficient terms. Only the threshold depends on \(N\); the
exponent used to choose the earlier source losses is fixed.
% qrh-collected-mask-cap-v1
For the joint source call, write the common puncture as
\(\mathfrak R_*L d_1d_2\), where \(\mathfrak R_*\) depends only
on the fixed character anchor, fixed bad-prime ideals and initial
puncture. The Newton assignment partitions one set
\(S\subseteq\mathcal P\), with \(\#S\le J_0\), so
\[
 d_1d_2=\prod_{\mathfrak p\in S}\mathfrak p,
 \qquad N(d_1d_2)\le Z^{J_0a}
 \quad\text{if }N\mathfrak p\le Z^a\quad(\mathfrak p\in\mathcal P).
\]
For any fixed \(\eta_*>0\), eventually
\(N\mathfrak R_*\le Z^{\eta_*}\). Consequently,
\[
 N(\mathfrak R_*L d_1d_2)\le Z^{4+J_0a+\eta_*}.
\]
The exponent is fixed before the ambient marked-slot count. Thus the
radical loss for this common puncture can be chosen within the fixed
source reserve. The moving missing ideal retains its own radical factor
in the weighted missing-ideal sum.

The full ideal harmonic bound now implies, for \(\delta>0\),
\[
 \sum_L\lVert\mu(L)\rVert\frac{(NL)^{\varepsilon_L}}{NL}
 \le H_{B_L,\delta}Z^{\delta+B_L\varepsilon_L}
 \qquad(\varepsilon_L\ge0).
\]
The constant and threshold precede the forcing pool and
\(\varepsilon_L\). Multiplication by the displayed paired-window
bound retains its complete allocation, profile and height costs.

Deleted marked scales remain in the exact raw
reduction factor. Apply the Newton and two-ideal expansions at this
completed positive norm, with their central factors and outside
character scalars in the same row Hilbert space.

\paragraph{Uniform admission at the moving puncture.}
For each extracted divisor pair \((d_1,d_2)\), use precisely
\(\mathfrak R d_1d_2\), the original keep predicate, the same outer
height and both rectangles shortened by the same divisors. The global
affine cap admits each positive shortening. Expanding the supplementary
holes in the common mask transports the completed ordinary bound from
its fixed character anchor. Uniformly over every surviving full marked
subset and every such shortening, for each fixed \(\delta>0\) the cost is
\[
 C_\delta\bigl(N\operatorname{rad}(\mathfrak R d_1d_2)\bigr)^\delta.
\]
The constant precedes the puncture, live characters and heights. The
previously established polynomial norm bound on the puncture absorbs
this factor in the original small power after \(\delta\) is chosen
inside the fixed reserve. The mesh precedes the ambient finite slot
type. Collision constants may depend on its later fixed cardinality;
finite profile controls and polynomial height degrees precede live rows
and heights, and the final large-height threshold follows those degrees.

% qrh-positive-radial-floor-v1
\paragraph{A fixed positive radial lower bound.}
Fix a node allowance \(e>0\). Let \(t,b\) be the original plain
and marked width bounds. Write \(Q=F-\delta-\rho\) in the common
geometry. For the actual retained row, the two radius allowances \(r,\eta\)
and frequency loss \(\omega\) give
\[
 Q=x_2=F-\delta-\rho,\qquad \rho=\omega+r-\eta,
 \qquad \operatorname{normValue}(z)\le U^Q.
\]
Choose \(r=\eta\), so that \(\rho=\omega\) and the existing frequency-loss
guard applies to this same radius. Let \(c\) be the sum of the three
common-ideal norm logarithms, and let \(f\) be the removed marked
width. Denote the annular, clipping and rounding allowance by
\(\ell\). For a nonzero positive corner, the supported clipped-scale
bound and nonnegative Newton norm logarithms give
\(0\le t-c+f+\ell\). Hence, using \(f\le b\),
\(c\le t+b+\ell\). The actual geometry also gives
\(0\le\delta\le c\). With \(0\le B\le1\) and the chosen guard
\[
 \ell+B\rho+2e\le BF-t-b,
\]
we obtain
\[
 BQ=BF-B\delta-B\rho
 \ge BF-c-B\rho\ge2e.
\]
This first implies \(Q>0\), and then \(BQ\le Q\) gives
\(Q\ge2e\). The other branch is an exactly zero positive corner.
Apply this argument separately to both corners. Thus the parent-source
lower endpoint \(M_{\min}=e\) is fixed from the node allowance
before the marked-slot cardinality; its eventual scale threshold may
depend on that cardinality.

% qrh-common-scale-grouping-v1
\paragraph{Common scales.}
For each positive corner put \(x_i(u)=\log_Z\max(1,X_i(u))\),
and let \(Q(u)\in[M_{\min},M_{\max}]\) be its radial exponent.
Partition the original rows by the three integer parts of
\((x_1(u),x_2(u),Q(u))/\omega\), with fixed \(\omega>0\).
In a nonempty cell let \(u_1,u_2,u_3\) be the three upper endpoints.
Then the common caps \(C_i=Z^{u_i}\) and
\(M=\min(M_{\max},u_3)\) satisfy, for every row in that cell,
\[
 X_i(u)\le C_i,\quad
 u_1+u_2\le x_1(u)+x_2(u)+2\omega,\quad
 M_{\min}\le M\le M_{\max},\quad Q(u)\le M\le Q(u)+\omega.
\]
The rounded affine admission includes the displayed \(2\omega\)
allowance. Increasing its radial exponent from \(Q(u)\) to \(M\)
preserves the downward face. Thus each cell has common admitted caps
and a common admitted radial scale. The polynomial geometry bounds
all three coordinates in a fixed interval, so the number of cells
is bounded independently of \(Z\) and of the original row family.
The weighted squared norm is exactly the sum over these cells;
their collection pays this fixed count and at most one additional
\(\omega\) in the radial exponent.

% qrh-attained-cell-weight-v1
In each nonempty cell, choose an original row attaining the largest
nonnegative weight in that cell, including its balanced first mass and
normalized allocation coefficient. The common radial and missing-ideal
upper exponents differ from those of this same row by at most the
scheduled cell widths. Apply the combined coefficient, missing-ideal,
radial and erasure price to this attaining row. The resulting bound
therefore retains the full scalar compensation. Empty cells contribute
zero. Sum the priced estimates over the finite cells and surviving
slot subsets, with their fixed cardinality factors.

% qrh-variable-restoration-family-v1
The shortening scales may vary with the original row \(u\).
Choose common positive caps \(C_i\) before all functions
\(0<Y_i(u)\le X_i(u)\le C_i\), with the same affine parent
length budgets. If the profile support is bounded above by \(b>0\),
put \(a=1/\max(1,b)\). A plain factor with scale below \(a\)
vanishes by its support. Otherwise replacing its scale by
\(\max(a,Y_i(u))\) leaves it unchanged. The norm comparison therefore
holds on every original row and retains its character, ideal mask,
numerator, slots and height. The common logarithmic windows cost
\[
 \prod_{i=1}^2\bigl(1+2\log(\max(1,C_i)/a)\bigr).
\]
These factors are uniform over the row-dependent shortenings.

For the forward two-ideal expansion take the union, over the original
finite rows, of their actual supported divisor-pair pools. Pad each
row's sum by its own membership indicator in this one union. Every
pair is good and inherits its norm cap from a witness row.
If \(b^2X_1(u)X_2(u)\le U^B\), the arbitrary-finite kernel mass
bound gives \(C_\delta U^\delta\) for this union. Finite-prime hole
deletion contributes its fixed mass. Weighted Cauchy now applies to
the whole row family at once: its restored squared norm is bounded
by the square of these masses, hence \(C_\delta U^{2\delta}\),
times the one shortened-family budget.
The polynomial caps on the \(C_i\), the divisor norms and the
puncture absorb the two logarithmic windows and these small powers
after their parameters are chosen within the fixed reserve.
% qrh-positive-corner-support-v1
For the centered row, first use
\(\lvert F_X-F_Y\rvert^2\le2\lvert F_X\rvert^2+2\lvert F_Y\rvert^2\).
For each positive corner restrict the original key family to the keys
where that corner is nonzero. Its actual supported terms then give
that corner's clipped-length bounds and earlier affine capacity budget.
Apply the one-corner ordinary source to the \(X\) family with
\((C_1^X,C_2^X)\), and separately to the \(Y\) family with
\((C_1^Y,C_2^Y)\). The zero terms vanish, and the original weights,
characters, masks, heights and singleton subsets are retained in both
families. The source constants are uniform before these restrictions.
Retain the corresponding explicit window budgets \(E_X,E_Y\).
With the common restoration mass \(\mathcal R\), weighted Cauchy gives
\[
 \sum_u\omega_u|F_X(u)-F_Y(u)|^2
 \le 2\mathcal R^2(E_X+E_Y).
\]
The equal products give the same central normalization for both
corners. Choose the larger constants, degrees and eventual thresholds,
and take the union of the finite profile controls. These choices are
uniform. Both rectangles retain their original outside scalars.


\hypertarget{sources:centered-fixed-inducing-terminals-retain-their-reference-saving}{%
\subsection{4. Centered fixed-inducing terminals retain their reference saving}\label{sources:centered-fixed-inducing-terminals-retain-their-reference-saving}}

Use the exact finite Newton expansion (CA8)--(CA9). Since \(0\le r_0\le R\), increase the eventual threshold to \(P\ge4R+4\). After common puncture extraction its shifted coefficients satisfy, for \(0\le j\le R\) and \(P\ge2(R+r_0+1)\), \[
 0\le c_j^{(r_0)}(P)\le 2^{j+1}j!P^{-j}.
\] There are at most \(P^j\) pool selections and at most \(2^j\) assignments of those selected divisors to the two ordinary factors. Since the product norm of an extracted divisor pair is at least one, the absolute weighted mass of the finite Newton kernel, at every positive exponent, is at most \[
 C_R=\sum_{j=0}^R2^{2j+1}j!.
 \tag{CA12}
\] This deliberately uniform bound suffices. Convolve that finite kernel with the exact forward two-ideal kernel (1)--(2). At exponent \(\sigma=1-\alpha=501/1000\), (6) and (CA12) bound its total divisor mass by \(O_R(U^\epsilon)\), uniformly over the common puncture.

Apply the ordinary equal-product centered estimate (MF3.78), interpolated with the absolute lattice bound as in the squarefree preprocessing section. For equal product scale \(T\) and minimum scale \(L>0\), it is \(O(\sqrt T L^{-\alpha})\), for \(\alpha=499/1000\), including subunit scales. A divisor pair of product norm \(d\) divides both rectangles by the same norms, so its centrally normalized contribution is \[
 d^{-1/2}\sqrt{T/d}\,(L/d)^{-\alpha}
   =\sqrt T L^{-\alpha}d^{-(1-\alpha)}.
 \tag{CA13}
\] Its extracted character is the same natural character, with an outside scalar of modulus at most one. The common puncture, Newton polynomial, profiles, heights and equal-product subtraction are identical for both rectangles. Sum (CA13) with the exact kernel mass at exponent .501. Taking the better absolute bound when the minimum scale is below one retains the relative terminal saving \[
 U^{-(499/1000)r+\epsilon},\tag{CA14}
\] where \(r\) is the nonnegative remaining reference exponent. With common width \(v\), this gives the saving \(\alpha(L_{\rm ref}-v)_+\) used in (G5). All finite profile controls and independent polynomial heights come from the same ordinary centered estimate before the divisor sum.

\hypertarget{sources:other-terminal-sectors-and-the-remaining-source-proof}{%
\subsection{5. Other terminal sectors and the remaining source proof}\label{sources:other-terminal-sectors-and-the-remaining-source-proof}}

The second Poisson true diagonal is already a sum of coefficient squares. On supported columns \texttt{P/(P-r)\textless{}=2}, so its diagonal estimate changes by at most a fixed factor four. The same direct coefficient bound applies to the coarse terminal absolute estimates. The artificial zero Gauss row retains its exact support; no zero is amplified into a nonzero contribution by (CA5).

\paragraph{An absolute envelope for Mellin inversion.}
Retain the actual finite second family, its corrected coefficients, both
ray characters and every residual ideal in the stored pool. Let \(l\)
index its labels, let \(w_l\) be its scalar weight, let \(o_l>0\) be
its physical outside factor, and let \(V_l\) be its input volume.
Write \(P_{00}=\|W_1\|_{0,0}\|W_2\|_{0,0}\).
The finite coefficient expansion and the lattice bound give
\[
 \sum_I |\widetilde b_l(I)|\le C_m P_{00}V_l.
\]
Here \(C_m\) depends on the fixed slot cardinality, coefficient and
support bounds. The correction has modulus at most two on the full
finite tuple pool, including its zero-column terms. If the physical
row norm is at most \(X\ge1\) and every stored ideal has norm at most
\(Y\), the full residual pool has cardinality at most \(128Y\).
Summing the nonnegative row majorant over the two ray characters gives
at most \(q^2D_0X\), where \(q\) is the fixed ray-character cardinality
and \(D_0\) is the diagonal control of the fixed ball profile.
Consequently the absolute second budget \(\mathcal A\) satisfies
\[
 \mathcal A\le128q^2D_0C_m^2P_{00}^2XY
       \sum_l |w_l|o_lV_l^2.
\]
This follows by the triangle inequality for the complete coefficient
sum, followed by the two finite counts. It retains the volume squared
for cancellation with the actual physical normalization.
For \(U>1\), if
\[
 X\le U^{B_r},\quad Y\le U^{B_p},\quad
 \#\{l\}\le C_lU^{B_l},\quad
 |w_l|o_lV_l^2\le C_wU^{B_w},\qquad B_r\ge0,
\]
then
\[
 \mathcal A\le128q^2D_0C_m^2C_lC_wP_{00}^2
                 U^{B_r+B_p+B_l+B_w}.
\]
The same bound holds for each normalized second sum, in either
orientation and at all real Mellin variables. All masks, natural
character zeros and original finite rows are retained. The stored
ideal pool has the exponent \(B_p=4\) established above.

\paragraph{Completing the Mellin integrals.}
Let \(H_0\ge1\) denote the external height factor and put
\(T=(1+|t|)(1+|u|)H_0\). Retain the original volume and the matched
bin factors \(E/H\) and \(EH\) in the two orientations. After their
scalar payments, write the bounded-range estimate as
\(C_s\mathcal F U^q\) for \(T\le U^\tau\), with \(\tau>0\),
and the corresponding absolute estimate as \(C_c\mathcal F U^A\).
Choose an integer \(J\ge0\) with \(\tau J\ge A-q\).
For \(T>U^\tau\), one has \(U^A\le T^JU^q\); the bounded range
has \(T^J\ge1\). Both ranges therefore satisfy
\[
 \text{budget}\le(C_s+C_c)\mathcal F T^JU^q.
\]
The finite Mellin moments of order \(J\) are controlled by finitely
many seminorms of the original smooth profiles. Minkowski's inequality
and Mellin inversion now give the physical norm bound, with a factor
\(H_0^J\) and the original weighted sum
\(\sum_k E(k)/(1+\operatorname{dyadicScale}(k_1))^d\).
The bin compensator stays paired with its own orientation until this
contraction. The polynomial exponents and \(J\) may depend on the fixed
slot type; they precede the moving profiles, rows and Mellin variables.

\paragraph{The MAIN zero and tail terms.}
Let \(V\) be the original volume and \(A_V=\log_U V\). For a node
\((B,\Lambda,h)\) with allowance \(e>0\), fix a width cell with
upper endpoint \(b\), rounding width \(\delta_w\), and
\(b-e-\delta_w\le\sum_iw_i\le b\). Put
\(t=B-\Lambda(b-e)\). The original affine admission gives
\[
 A_V\le t+b+\Lambda\delta_w.
\]
The parent-budget certificate, \(B_0\le1\), and
\(\delta_w\le e/(16(\Lambda+1))\) give
\[
 t+b+4e<B_0(1-6a)\le1-6a,
 \qquad A_V+6a+3e\le1.
\]
The actual MAIN radius has exponent \(R+6a+\omega\), where
\(\omega\) is the frequency and support allowance. The balanced-side
and common-radical identity pays its base exponent by \(A_V-1\).
The reciprocal amplifier-pool cardinality contributes \(-a\); the
annular band count is at most seven. Both allocation-cardinality
payments are retained. After the remaining small counting losses
\(\ell\), the zero-term exponent is therefore at most
\[
 A_V-1+5a+\omega+\ell\le-a
 \qquad\text{when }\omega+\ell\le3e.
\]
The tail with its original saving \(s>0\) has exponent at most
\(-a-s\). The previously displayed profile and height factors remain
outside these powers of \(U\).

All MAIN radius, raw/natural width, central normalization and common allocation formulas are unchanged. The resulting source certificate checks MAIN HIGH, MAIN LOW, RAW admission, fixed and tiny-frequency sectors, the squarefree centered terminal, and each shortened-factor fallback. The finite Newton multiplier and the squarefree restriction are expanded only at the explicit completed norm stage described above; fixed-inducing terminals retain equal-product cancellation in (CA13).

\hypertarget{sources:chapter-27.-restricted-high-and-every-exceptional-sector}{%
\section{Chapter 27. Restricted HIGH and every exceptional sector}\label{sources:chapter-27.-restricted-high-and-every-exceptional-sector}}

\hypertarget{sources:restricted-high-counting-root-completion-and-exceptional-sectors}{%
\section{Restricted HIGH counting, root completion, and exceptional sectors}\label{sources:restricted-high-counting-root-completion-and-exceptional-sectors}}

Use the density interval

\[
                    \frac12\le\nu\le\frac34.             \tag{LD1}
\]

Use inverse Theorem 7.3, plain Theorem 7.4, and the common-zero witness Lemma 14.2. The ordinary and complete marked squarefree projectors give exactly the four first-common types (10).

\hypertarget{sources:coefficients-and-the-exact-witness-calculation}{%
\subsection{1. Coefficients and the exact witness calculation}\label{sources:coefficients-and-the-exact-witness-calculation}}

Write

\[
 d=5-2\nu,\quad w=5-6\nu,\quad
 \rho=1-\frac{10\nu}{3d},\quad \phi=\frac{\nu w}{d},
\]

\[
 A_\nu=\frac{14\nu w}{115-54\nu},\quad
 B_\nu=\frac{2\nu w^2}{3d(115-54\nu)},\quad
 C_\nu=\frac{16\nu w^2}{3d(185-138\nu)}.
\]

All denominators are positive on (LD1). Define

\[
 G_\nu(S,Q)=\min\{\phi S,A_\nu S+B_\nu Q,C_\nu Q\}.
                                                               \tag{LD2}
\]

Normalize Q=1, first take amplitude parameter p=1/2, and use exactly the inverse and ordinary witnesses of Lemma 14.2. Their available marked widths are

\[
 h_I=\min(S,(1-r)/2),\qquad
 h_P=\min(S,2(1-2m)/9).
\]

The count bounds, when invoked, are

\[
 E\le1-\nu\min(r+S,(1+r)/2),\quad
 E\le1-\nu\min(2m+S,(14m+2)/9).
\]

For a target xi put

\[
 r_*=\max(\xi-S,2\xi-1),\quad
 y_*=\max(\xi-S,(9\xi-2)/7),\quad
 t_*=r_*+y_*/2.
\]

The equation

\[
 1-\nu\xi=\frac16+(5/6-\nu)t_*                 \tag{LD3}
\]

has the following three branches:

\[
 \xi=\begin{cases}
 (10+3wS)/(15-6\nu),&0\le S\le s_1,\\
 (70+2w+14wS)/(115-54\nu),&s_1\le S\le s_2,\\
 (70+16w)/(185-138\nu),&S\ge s_2,
 \end{cases}
\]

\[
 s_1=\frac{2w}{3(45-26\nu)},\qquad
 s_2=\frac{7w}{185-138\nu}.                         \tag{LD4}
\]

These are rational identities, not new analytic inputs. Their joins agree, 0\textless=s1\textless=s2, and the plain cap becomes active before the inverse cap. Substitution gives

\[
 \rho-(1-\nu\xi)=G_\nu(S,1).                         \tag{LD5}
\]

\hypertarget{sources:every-witness-admission-remains-strict}{%
\subsection{2. Every witness admission remains strict}\label{sources:every-witness-admission-remains-strict}}

On (LD1), the inverse second gate is checked as follows:

\[
 \xi\ge\frac{10}{3(5-2\nu)}\ge\frac56,
 \qquad r_*\ge\frac{22}{29}>\frac34.              \tag{LD6}
\]

For completeness, the three expressions for \(\xi\) in (LD4)
give the last bound directly. If \(0\le S\le7/58\), clearing their
positive denominators gives \(\xi_i-S\ge22/29\) for each branch.
If \(S\ge7/58\), each branch satisfies \(\xi_i\ge51/58\).
Taking their minimum and then
\(r_*=\max(\xi-S,2\xi-1)\) proves (LD6) on the whole density interval.
The inverse moment is invoked only when \(r\ge r_*\). Its second
capacity therefore satisfies

\[
 2r+8h_I\le4-2r,\qquad
 3-(2r+8h_I)\ge\frac{15}{29}.                     \tag{LD7}
\]

The first inverse gate and ordinary marked gate receive the same fixed width erosions described after (LD8). The unused second inverse margin (LD7) uniformly absorbs the witness enclosure error.

The largest xi is its saturated value

\[
 \xi_{\rm sat}=\frac{150-96\nu}{185-138\nu}<1,
\]

because 35-42nu\textgreater0 on (LD1). Hence (LD3), w\textgreater0, and (LD5) imply

\[
                   1<t_*\le\frac5{5-2\nu}\le\frac{10}{7}<\frac32.
                                                               \tag{LD8}
\]

This is an admitted cutoff for the same-zero detector. If the inverse witness is longer than one, its original unmarked long bound is monotone in r since 5/6-nu\textgreater0, and r\textless=t* pays it. If the ordinary witness exceeds one half only by its admitted enclosure, its unmarked bound suffices because xi\textless1. Thus no witness outside a moment's domain has been used.

The passage from p=1/2 to 0\textless=p\textless=1/2 is unchanged: subtract nu\emph{(1-2p)}S from the count exponent, and use 2p\emph{h+(1-2p)}S\textgreater=h for every selected width h\textless=S. All marked subsets remain whole original slots with the canonical pools and coefficients. Consequently the natural common-zero count is

\[
             \rho Q+\nu S-G_\nu(S,Q)+\varepsilon.       \tag{LD9}
\]

\hypertarget{sources:uniform-coefficient-bounds-and-full-sixth-root-completion}{%
\subsection{3. Uniform coefficient bounds and full sixth-root completion}\label{sources:uniform-coefficient-bounds-and-full-sixth-root-completion}}

Direct differentiation shows rho decreases on (LD1), so

\[
 \frac27\le\rho\le\frac7{12},\qquad
                 0\le\phi<\frac{13}{50}.             \tag{LD10}
\]

The second bound holds for every real nu with positive d: after clearing d, its strict difference is the positive quadratic 300nu\^{}2-276nu+65, whose discriminant is -1824. The bound applies to the whole interval (LD1). The natural plain capacity remains at most 9/2; the density \ensuremath{\nu} does not replace that capacity.

Put \(x=\nu-1/2\in[0,1/4]\) and
\(D=3(5-2\nu)(185-138\nu)>0\). The exact identity
\[
 D-40\bigl(16\nu(5-6\nu)^2\bigr)
 =112+2495x+1092x(1/4-x)+23040x^2(1/4-x)>0
\]
proves the uniform coefficient bound
\[
                    0<C_\nu<\frac1{40}.                 \tag{LD11}
\]
Clearing positive denominators also gives
\(0<A_\nu<\phi\) and \(0<B_\nu<C_\nu\).
Thus \(G\) is monotone, \(\phi\)-Lipschitz in marked width and
\(C_\nu\)-Lipschitz in row width.

For the two root sums use \(y=3/4-\nu\ge0\). Direct expansion gives
\[
\begin{aligned}
175D\bigl(6(\rho-C_\nu)-279/175\bigr)
 &=59283/4+1191798y+1734588y^2+604800y^3,\\
350D\bigl((\rho-1/6)/2-C_\nu-83/2100\bigr)
 &=19761/4+193516y+233196y^2+201600y^3.
\end{aligned}
\]
Both right sides are positive. Hence the small-root series exponent and
large-root tail reserve satisfy
\[
                   6(\rho-C_\nu)>\frac{279}{175}>1,       \tag{LD12}
\]
\[
              (\rho-1/6)/2-C_\nu>\frac{83}{2100}>0.       \tag{LD13}
\]
These estimates hold uniformly on the full interval (LD1).
Here is the full root summation, with the same masks and independent
heights. Write each actual row uniquely as \(\zeta uB^6\), retaining
its six-unit label and sixth-free core, and put \(n=NB\).
At total raw width \(D\in[D_{\min},D_{\max}]\),
\(D_{\min}>0\), its residual core width is
\(q=D-6\log_U n\). For \(n\le U^{D/12}\),
\(q\ge D/2\). Choose the core mesh at most
\(\mu D_{\min}/2\) before roots or heights.
Since \((U^q)^{w_j/q}=U^{w_j}\), the actual slot scales and common
profiles are unchanged. The fixed positive core lower width permits
one common height exponent, finite seminorm set and eventual threshold.

The exact deletion of root and moving-hole primes from a slot costs \(O(P_j^{-1/2}\log U)\): subtract precisely its terms with a prime dividing the hole. After the previously assigned tiny-slot cutoff this tends to zero uniformly for polynomially bounded holes. The core character and relative coefficients remain the same; only their natural zero factors change. The finite positive maximum-product comparison cost is at most \(4^{|J|}\), a constant after the slot system is fixed.

Let \(g\) be the row-width Lipschitz bound of the gain. Then \(G(L,q)\ge G(L,D)-6g\log_U n\), so a small-root contribution is at most \[
 U^{\rho D+\nu L-G(L,D)+\epsilon}n^{-6(\rho-g)}.
 \tag{RT1}
\] The ideal sum converges when \(6(\rho-g)>1\), with no additional root cardinality. For \(n>U^{D/12}\), use the unmarked core count and the pointwise marked envelope. Partial summation of its ideal root series saves \((\rho-1/6)D/2\). Since \(G(L,D)\le gD\), this pays the full gain whenever \((\rho-1/6)/2-g>0\). Empty subunit residual windows contribute zero. The fixed finite nonprincipal exceptional cores cost their marked envelope and are covered by \(U^{(\rho-g)D/2}\ge1\).

For the three-facet gain take \(g=C_\nu\), using (LD12)--(LD13). For the later six-facet gain take \(g=C_{\max}\), using (HF13)--(HF14). Fixed-inducing cores retain their separate count \(U^{Q/6+L}\); tiny nonprincipal cores retain \(U^{Q+L}\). Their complete payments are (LD15)--(LD17). This proves full-gain completion for each gain used in the source graph.

\hypertarget{sources:the-restricted-first-common-cone-is-individually-favorable}{%
\subsection{4. The restricted first-common cone is individually favorable}\label{sources:the-restricted-first-common-cone-is-individually-favorable}}

Let A,b,c,D be the total logarithmic widths of genuine common primes of types (1,1),(1,2),(2,1),(2,2), respectively. The total first lengths are

\[
 C=A+b+2c+2D,\qquad D_c=A+2b+c+2D.
\]

The exact first-local rules give R=b+c, L\_\allowbreak{}C=2c, R0\_\allowbreak{}C=b, Sigma=D, and both B positive parts vanish. The complete common form before the density charge is

\[
 V_C=(1-\rho)(C-D_c)_+-(1+\rho)R+\rho L_C+R_{0,C},
\]

one has the global identity

\[
 V_C-\Sigma=-\min(b,c)-\rho(b-c)_+-D\le0.              \tag{LD14}
\]

This holds for every 0\textless=rho\textless=1. It does not require primewise ordering. The reverse side has the same identity with b,c exchanged.

For an ordinary HIGH side, conserved second-extraction width pays the following exact Lipschitz payment. If \(L+c+b_2\ge S\) and \(Q+2c+2b_2\ge F\), monotonicity, clipping, and the width Lipschitz constants give \[
 \nu b_2+G_\nu(L,Q)
 \ge G_\nu((S-c)_+,(F-2c)_+)
 \ge G_\nu(S,F)-(\phi+2C_\nu)c.
 \tag{LD14a}
\] To prove the first inequality, put back second common width \(b_2\): its possible gain cost is at most \((\phi+2C_\nu)b_2\le\nu b_2\). The second inequality puts back the first common width \(c\) and uses the same constants. This argument includes all clipped regimes. The bound is

\[
 \phi+2C_\nu<\frac{13}{50}+\frac9{100}=\frac7{20}<\nu.
\]

Equivalently use eta=nu-phi-9/100\textgreater1/10. Equations (LD14) and this inequality make the residual ordinary HIGH common excess individually nonpositive after the three no-common scalar gates H+G\_\allowbreak{}nu(S,F)\textgreater0 have been certified.

\hypertarget{sources:fixed-inducing-and-tiny-radius-raw-sectors-remain-present-and-paid}{%
\subsection{5. Fixed-inducing and tiny-radius raw sectors remain present and paid}\label{sources:fixed-inducing-and-tiny-radius-raw-sectors-remain-present-and-paid}}

For a fixed-inducing core replace the ordinary core count by its actual exponent Q/6+L. The same inner allowance (140) and common payment (RC2) then give the excess

\[
 -H+V_C-\Sigma-\nu(C+b_2)
       -(\rho-1/6)Q+(1-\nu)L.                        \tag{LD15}
\]

Use (LD14), the original geometric relation 2(C+b2)+Q\textgreater=F, and L\textless=s. With u=C+b2, the fixed-inducing reserve is at least

\[
 H+\nu u+(\rho-1/6)Q-(1-\nu)s.
\]

It is paid by the two sufficient scalar gates

\[
 H+\frac{\nu F}{2}-(1-\nu)s>0,\qquad
 H+(\rho-1/6)F-(1-\nu)s>0.                        \tag{LD16}
\]

The count Q/6+L retains the actual fixed cores, their units, and every sixth-power repetition. In a fixed nonprincipal tiny raw ball Q\textless Qmin, use the original absolute count Q+L. Its sufficient gate is

\[
 H+\frac{\nu(F-Q_{\min})}{2}
           -(1-\rho)Q_{\min}-(1-\nu)s>0.           \tag{LD17}
\]

Choose Qmin\textgreater0 before all live labels. On the complementary positive radius, (LD9)--(LD13) apply uniformly. If a LOW partner can have positive natural excess Delta, subtract max(Delta,0) from each special reserve; one cannot replace that deduction by bare positivity. If the admitted RAW route has already made every LOW individually nonpositive, no such positive partner remains. Ordinary HIGH partners are individually nonpositive by Section 4. Two special sectors are each paid by their own scalar gates.

\hypertarget{sources:chapter-28.-the-global-ordinary-source-theorem}{%
\section{Chapter 28. The global ordinary source theorem}\label{sources:chapter-28.-the-global-ordinary-source-theorem}}

\hypertarget{sources:a-global-affine-source-family-and-a-finite-acyclic-construction}{%
\section{A global affine source family and a finite acyclic construction}\label{sources:a-global-affine-source-family-and-a-finite-acyclic-construction}}

\hypertarget{sources:statement-of-the-inductive-family}{%
\subsection{1. Statement of the inductive family}\label{sources:statement-of-the-inductive-family}}

Retain the literal row character, all natural zero sets, fixed-inducing exclusions, ray factors, two ordinary ideal sums, and canonical marked prime polynomials (158)--(160). Canonical marked prime pools may overlap. All heights and smooth profiles have the same finite seminorm conventions of Section 1.3 and (MF3.40).

Write \texttt{G(B,Lambda,h)} for the following \textbf{ordinary} initial moment:

\[
 \mathcal E\ll U^{1+\epsilon}\mathcal H_J
 \quad\text{whenever}\quad
 m_1,m_2\ge0,\quad 0\le s\le h,\quad
 m_1+m_2+\Lambda s\le B.
\tag{G1}
\]

There are no individual lower or upper plain edges in this family. Empty and bounded subunit annuli are treated with the inherited clipping convention. Initial actual nonempty profiles under (G1) have polynomial norm, since the construction keeps \texttt{B\textless{}=1}.

The first input is the already proved natural theorem \texttt{G(1,K,h0)}, where \texttt{K=6(2\ theta\_\allowbreak{}I-1)} and any required fixed \texttt{h0} is permitted by its natural capacity inequality. Every later \texttt{G} will use one strictly earlier \texttt{G}, together with the natural theorem (MF3.40) and the inherited marked HIGH count. Its exact source certificate records the earlier parent ID, its rational affine data and its entire width interval. The construction is a finite directed acyclic graph.

\hypertarget{sources:restricted-preprocessing-and-global-completion}{%
\subsection{2. Restricted preprocessing and global completion}\label{sources:restricted-preprocessing-and-global-completion}}

First prove the current moment with \(\mu^2(l_1l_2)\) on the ordinary product and the retained squarefree projector on the complete marked product. The finite collision decomposition (10a) and slot-cardinality induction then recover the unrestricted marked source. The exact forward and inverse kernels, central masses, and extraction law are proved in (1)--(9) of the squarefree preprocessing section.

The inverse kernel expresses the desired ordinary row as a sum of restricted rows, with absolute central mass \texttt{U\^{}epsilon}. Each of its two plain scales is only shortened. Because the new statement (G1) contains every shorter pair under the total cap, every such row is admitted. Its extracted character is a contracting scalar in the same positive row Hilbert space. This transfers the completed restricted theorem to the ordinary theorem (G1); no pointwise comparison between uncentered coefficient norms is used.

The forward restriction is kept through the first common extraction, amplification and the second row transform. It is restored only at a completed non-fixed-inducing norm to which an earlier moment applies. Fixed-inducing terminals retain the whole restricted centered coefficient. The exact extraction law preserves the restriction plus a common puncture; a non-squarefree extracted ordinary product contributes zero before positivity. Thus all genuine first common multiplicities belong to \texttt{(1,1),(1,2),(2,1),(2,2)}.

\hypertarget{sources:initial-short-plains-and-an-adjustable-reference}{%
\subsection{3. Initial short plains and an adjustable reference}\label{sources:initial-short-plains-and-an-adjustable-reference}}

For a source width cell with upper width \texttt{b}, choose one fixed reference

\[
 L(b)=\frac{49}{100}-Kb>0.
\tag{G2}
\]

If either original plain is shorter than \texttt{L(b)}, restore the restricted row by its forward kernel and use the natural reflection (MF3.44) on the other ordinary factor. The first factor remains no longer than \texttt{L(b)}. After reflection, whole-factor conjugation and the inherited profile control, the other factor is no longer than \texttt{1/2+o(1)}. The natural capacity is therefore at most

\[
 L(b)+\frac12+Kb+o(1)=\frac{99}{100}+o(1)<1.
\tag{G3}
\]

The same reasoning applies to every shifted forward-kernel summand: the first factor is only shorter and the reflected bound for the second is still one half. All character zeros and central factors are those of natural reflection (MF3.44). One reflects only after forming an ordinary factorized product; the squarefree projector itself is not treated as a separable smooth coefficient.

It remains to consider original plains at least \texttt{L(b)}. If their actual total is \texttt{Pact}, use the centered reference rectangle

\[
 Y_1=U^{L(b)},\qquad Y_2=U^{Pact-L(b)}.
\tag{G4}
\]

Here \texttt{Pact\textgreater{}=2L(b)}, so both reference lengths are at least \texttt{L(b)} and their product agrees exactly with that of the original rectangle. Its initial energy is bounded by the same restoration/reflection argument (G3). The first exact source identity is then applied to the entire restricted centered difference.

An earlier global theorem (G1), when called at a completed transformed RAW norm, admits either rectangle directly from its total-cap inequality.

\hypertarget{sources:elementary-centered-terminal-bound}{%
\subsection{4. Elementary centered terminal bound}\label{sources:elementary-centered-terminal-bound}}

Apply the equal-product ordinary centered estimate and its exact divisor transport proved in the preprocessing section, with \(\alpha=499/1000\) and mass exponent \(1-\alpha=501/1000\). Both rectangles retain the same character, puncture and Newton multiplier. The common leading term cancels before absolute values; all finite profile seminorms and independent heights are kept. This gives the relative saving \(U^{-\alpha r+\epsilon}\) at every fixed-inducing terminal.

The terminal expression with common width \texttt{v} is bounded by

\[
 A-\frac56-\frac23v-\alpha(L(b)-v)_+.
\]

Since \texttt{0\textless{}alpha\textless{}2/3}, its maximum over \texttt{v\textgreater{}=0} is at \texttt{v=0}. Including the amplifier and support costs in (131)--(141), the sufficient centered terminal gate is therefore

\[
 \boxed{\quad \frac56+\alpha L(b)-T-9a>0,\quad T=t+b.\quad}
\tag{G5}
\]

The final schedule checks this width-dependent gate directly at both endpoints of every source interval. The divisor transport uses the absolutely convergent mass at exponent 501/1000. All fixed-character leading terms are retained until the centered subtraction is made.

\hypertarget{sources:exact-common-geometry}{%
\subsection{5. Exact common geometry}\label{sources:exact-common-geometry}}

Let \texttt{A,B12,C21,D} be the nonnegative log-widths of the four permitted first common types. Then

\[
 \begin{aligned}
 c&=A+B_{12}+2C_{21}+2D,&
 d&=A+2B_{12}+C_{21}+2D,\\
 R&=B_{12}+C_{21},&
 L_c&=2C_{21},&R_{0,c}&=B_{12},\\
 B_c&=B_d=0,&\Sigma&=D.
 \end{aligned}
\tag{G6}
\]

Put \texttt{delta=(d-c)\_+=(B12-C21)\_+}. The actual reconstructed RAW norm is

\[
 Q=F-\delta,\qquad F=1-6a.
\tag{G7}
\]

The reverse orientation has the same formulas with \texttt{B12,C21} exchanged. In particular \texttt{delta\textless{}=B12}, \texttt{delta\textless{}=c}, and \texttt{2delta\textless{}=d}. The complete common ideal divides an actual original column, so \texttt{d\textless{}=T+o(1)}. Thus

\[
 Q\ge F-T/2-o(1).
\tag{G8}
\]

The exact certificate includes \texttt{F-T/2\textgreater{}0} with reserve. Earlier source moments are therefore called at a scale \texttt{U\textquotesingle{}=U\^{}Q} with a uniformly positive exponent. Their own natural reflection and reference arguments are expressed in that scale. Tiny individual actual rows inside this ball do not change the ambient source scale. The separate small-core HIGH sectors are still checked.

% qrh-first-centered-payment-v1
\paragraph{The first centered payment.}
Put \(A_0=\log_Z V\) for the original parent volume,
\(p=\log_Z N\operatorname{rad}(\mathfrak C)\), and
\(L=49/100-Kb\). Let \(e\ge0\) be the first projector exponent.
With the common widths \(c,d,R\) of (G6), define
\[
 \Pi=\frac R3-\frac e6-\frac{2c}3-\frac{4d}3
       +\frac56(d-c)_+.
\]
The four nonnegative type widths in (G6) give
\(\Pi+p\le-2c/3\), in either orientation.
The actual first MAIN radius has logarithm
\[
 2A_0-1+R+e-c-d+(d-c)_++6a+\omega,
\]
where \(\omega\) is the retained frequency and radius reserve.
After the original balanced normalization, one prime average and
all explicit remaining small losses \(\ell\), its centered exponent is
\[
 A_0-\frac56+4a+\frac56\omega+\ell+\Pi
       -2\alpha(L-c)_+.
\]
Here the clipped reference length is exactly
\(\max(1,Z^L/N\mathfrak C)=Z^{(L-c)_+}\), including subunit ratios.
For \(0\le\alpha\le2/3\) and \(c\ge0\),
\[
 -\frac{2c}3-2\alpha(L-c)_+\le-\alpha L.
\]
Write \(h_H\) for the prescribed HIGH source gain. Choose the
small errors so that \(A_0\le T+\delta\) and
\(\delta+5\omega/6+\ell\le a\). The exact schedule checks the
charged form of (G5),
\[
 T+9a+h_H<\frac56+\alpha L.
\]
Consequently the displayed exponent is at most
\(-h_H-4a-p\). The remaining first-key and projector sum retains
its original balanced scalar and is bounded by the ideal harmonic
sum against \(Z^{-p}\), with an arbitrarily small scale loss.
The four amplification widths left above absorb that loss and the
original first allocation, band and fixed-cell factors.

\hypertarget{sources:pure-raw-low-with-both-floors-and-a-finite-upper-slot-cap}{%
\subsection{6. Pure RAW LOW with both floors and a finite upper slot cap}\label{sources:pure-raw-low-with-both-floors-and-a-finite-upper-slot-cap}}

Let the strictly earlier theorem be \texttt{G(B0,Lambda0,h0)}, with \texttt{B0\textless{}=1} and \texttt{Lambda0\textgreater{}1}. Fix a LOW threshold \texttt{q}, a positive small reserve \texttt{r}, and choose

\[
 a=\frac{q\{b+(t-B_0)/\Lambda_0\}+r}
              {1-6qB_0/\Lambda_0},\qquad
 g=a-q\{b+(t-B_0F)/\Lambda_0\}=r.
\tag{G9}
\]

The denominator, amplifier range and padding reserves are explicit certificate gates. Let \texttt{f} be the first frozen marked width, so that \texttt{0\textless{}=f\textless{}=min(c,b)}. The first quotient caps are

\[
 P_1=t-c+f,\quad s_1=b-f,\quad
 z_0=\frac{B_0F-t}{\Lambda_0},\quad
 z=\frac{B_0Q-P_1}{\Lambda_0}
   =z_0+\frac{c-B_0\delta-f}{\Lambda_0}.
\tag{G10}
\]

Require the global scalar guards

\[
 B_0F-t-b>0,\qquad z_0\le h_0F,
 \qquad qh_0\le1.
\tag{G11}
\]

The first implies \texttt{z\textgreater{}0} in every genuine common geometry, since \texttt{c-B0\ delta\textgreater{}=0} and \texttt{f\textless{}=b}. Retain LOW slots only up to \texttt{min(z,h0\ Q)}; subsequent common or artificial extractions only shorten the two plains and remove slots. The earlier global theorem applies to every resulting positive completed row norm. Greedy mesh overshoot is charged to the fixed small losses.

For the actual joint row (168), the first allowance (MF3.55), the second allowance (140), and the shared first payment (RC2) give the complete RAW excess before erasure: \[
 -a-2R-E-B_c+L_c+R_{0,c}+g_2-t_2-b_2-\Sigma.
 \tag{RC3}
\] Here \(E\ge0\) is the inactive projector width; the second extraction has \(b_2-g_2+t_2\ge0\), as follows from its divisor and central normalizers in (138)--(141). The source theorem is applied to the whole row (168), with its fixed ray character and puncture, after regrouping actual source labels by that row. The multiplicity bound in the literal row section pays the missing first radical. Using (G6), (RC3) is at most \(-a-B_{12}-D\). The entire positive-part erasure cost is therefore bounded by

\[
 E_c\le-a-B_{12}-D
    +q\max\{0,b-f-z,b-f-h_0Q\}+o(1).
\tag{G12}
\]

This maximum retains both the zero-erasure floor and the prior's upper marked-width cap. Its three branches are paid separately:

\[
 \begin{aligned}
 E_{c,0}&\le-a,\\
 E_{c,z}&=-g-B_{12}-D-\frac q{\Lambda_0}(c-B_0\delta)
             -q(1-1/\Lambda_0)f\le-g,\\
 E_{c,h}&=-g+q(z_0-h_0F)+qh_0\delta-B_{12}-D-qf
             \le-g.
 \end{aligned}
\tag{G13}
\]

For the last inequality use \texttt{delta\textless{}=B12} and (G11). Consequently each orientation has the individual LOW reserve \texttt{min(a,g)}, before the scheduled small losses. The third branch handles a normalized slot width above \texttt{h0} in the given common geometry.

The three branches prove every earlier-face call on its full global domain, including shortening and cap branches.

\hypertarget{sources:high-and-the-remaining-gates}{%
\subsection{7. HIGH and the remaining gates}\label{sources:high-and-the-remaining-gates}}

Use the three-branch marked HIGH gain (LD2) with

\[
 \rho=1-\frac{10\nu}{3(5-2\nu)},\quad
 \phi=\frac{\nu(5-6\nu)}{5-2\nu},
\]

\[
 A_\nu=\frac{14\nu(5-6\nu)}{115-54\nu},\quad
 D_\nu=\frac{2\nu(5-6\nu)^2}
              {3(5-2\nu)(115-54\nu)},\quad
 C_\nu=\frac{16\nu(5-6\nu)^2}
              {3(5-2\nu)(185-138\nu)}.
\tag{G14}
\]

The inverse admissions and full-root completion on \texttt{1/2\textless{}=nu\textless{}=3/4} are proved in (LD3)--(LD13) and the root summation following them. In this restricted first common class the balanced common expression is

\[
 V_c-\Sigma=-\min(B_{12},C_{21})
             -\rho(B_{12}-C_{21})_+-D\le0.
\tag{G15}
\]

It therefore pays \texttt{nu*c} individually before marked-width losses. The exact certificate checks \texttt{nu\textgreater{}phi+2Cnu} on the full density interval and the three regular gates, for \texttt{H=1-rho-nu*T-a(5-6rho)},

\[
 H+\phi S>0,\quad H+A_\nu S+D_\nu F>0,
 \quad H+C_\nu F>0.
\tag{G16}
\]

The fixed-inducing and tiny-core HIGH sectors are retained with the individual common rate \texttt{nu}:

\[
 \begin{gathered}
 H+\nu F/2-(1-\nu)b>0,\\
 H+(\rho-1/6)F-(1-\nu)b>0,\\
 H+\nu(F-r_0)/2-(1-\rho)r_0-(1-\nu)b>0,
 \qquad r_0=1/50.
 \end{gathered}
\tag{G17}
\]

Here \texttt{F,S} in a certificate include its explicit downward padding. Other true-diagonal, tiny-frequency and reference-extraction guards are \texttt{1-T-5a}, \texttt{1-r0-T-4a}, \texttt{1-r0-T-(9/2)a}, and \texttt{L(b)-7a}, each with its stated positive fee. The corrected amplifier has only MAIN, by the exact corrected-amplification identity (CA6).

All quantities are affine in the width-cell endpoint \texttt{b} when \texttt{q} and the two source faces are fixed. After multiplying (G16)--(G17) by their positive density denominators, they are polynomials of degree at most three in \texttt{nu}. Exact Bernstein positivity at both width endpoints therefore proves the full continuous rectangle of cell widths and densities. Subdivision in density, when needed, remains an exact finite certificate.

\hypertarget{sources:acyclic-inheritance-and-uniform-completion}{%
\subsection{8. Acyclic inheritance and uniform completion}\label{sources:acyclic-inheritance-and-uniform-completion}}

A new affine face may be weaker than its selected parent below their exact rational crossing. On that lower interval use the parent theorem directly; its cap must cover the entire inheritance interval. Above the crossing, certify the source gates at both width endpoints. This proves the new face on its complete interval \texttt{{[}0,h{]}}, including every smaller marked width. Each parent is an earlier node, never a proposed current or later face.

Choose all arithmetic-independent width meshes, positive amplifier lower bounds, row/radius losses, fixed-character contour orders, profile orders and height exponents over this finite graph before the moving slot labels and target height. The common mesh bounds every marked slot width below each fixed amplifier lower bound, before the finite slot type is chosen. The original Fourier-window, ray, unit, puncture, first allocation and positive row norm conventions are those of (MF3.40), (120)--(141), and (158)--(168). The all-height completion is applied to the original marked initial
energy. Write \(N\) for the ambient slot cardinality,
\(H=1+|t|+\mathrm{height}\ge1\), and
\(\Delta_S=\operatorname{diagonalControl}(\Phi)\,
\operatorname{control}_S(p)^2\). The elementary absolute bound gives
\[
 E\le C_c\Delta_{S_c}H^{j_c}U^{4+2N\mu},
\]
where \(\mu\) is the fixed marked-width mesh. The bounded-height
physical estimate and reference reduction give
\[
 H\le U^\tau\quad\Longrightarrow\quad
 E\le C_s\Delta_{S_s}H^{j_s}U^{1+\varepsilon}
\]
for some \(\tau>0\). Choose an integer \(n\ge0\) with
\(\tau n\ge3+2N\mu-\varepsilon\), and put
\(J=\max(j_s,j_c+n)\), \(S=S_s\cup S_c\).
If \(H>U^\tau\), then
\(U^{4+2N\mu}\le H^nU^{1+\varepsilon}\). Thus both cases give
\[
 E\le(C_s+C_c)\Delta_S H^J U^{1+\varepsilon}.
\]
The source mesh and frequency allowance are fixed before the slot type.
After that type and the arithmetic anchors are fixed, the bounded-height
cutoff, finite controls, degrees and constants precede the eventual
scale threshold and all moving profiles, characters and heights.
Finite minima of the positive cutoffs handle all original slot subsets
and width cells. The energy, its puncture, keep predicate, reference
rectangle and radial scale remain the original ones.

The finite rational gates below and the preceding analytic arguments prove each new admitted initial moment.

The numerical certificate's printed width, radius and analytic accounting caps are each at most \(10^{-12}\). They are upper allowances for the actual losses chosen below. First take the minimum of the finitely many strict rational reserves, after normalizing the cleared density denominators. Then choose the actual support buffers, Fourier losses, greedy slot overshoot and asymptotic error exponents small enough that their total contribution is below those allowances and that minimum reserve. Such a choice exists by continuity and finiteness. In a width cell, \(q\), \(a\) and the reference length are chosen from its fixed upper endpoint, before any moving prime labels or heights.

The centering parameter used by this theorem and its checker is exactly \(\alpha=499/1000\).

\hypertarget{sources:chapter-29.-the-complete-finite-rational-source-schedule}{%
\section{Chapter 29. The complete finite rational source schedule}\label{sources:chapter-29.-the-complete-finite-rational-source-schedule}}

\hypertarget{sources:exact-finite-source-schedule}{%
\section{Exact finite source schedule}\label{sources:exact-finite-source-schedule}}

\hypertarget{sources:final-admitted-source-moments}{%
\subsection{1. Final admitted source moments}\label{sources:final-admitted-source-moments}}

The following are proved ordinary initial source moments. They retain the literal row character, natural zero sets, fixed-inducing exclusions, all-height profile conventions and marked polynomial conventions of the literal row (160)--(161). Canonical marked prime pools may overlap. The source bound is

\[
 \mathcal E\ll U^{1+\epsilon}\mathcal H_J
 \quad\text{if}\quad
 m_1,m_2\ge0,\quad 0\le s\le h,\quad
 m_1+m_2+\Lambda s\le B.
\tag{S1}
\]

In particular every shorter plain pair and every smaller marked width is included. No lower marked-width edge or individual plain floor is part of this source contract.

\begin{landscape}
\begingroup % exact-table-layout
\fontsize{7.5}{9.5}\selectfont\setlength{\tabcolsep}{2pt}\renewcommand{\arraystretch}{1.18}
\begin{longtable}{@{}
 >{\raggedright\arraybackslash}p{42mm}
 >{\raggedright\arraybackslash}p{42mm}
 >{\raggedright\arraybackslash}p{34mm}
 >{\raggedright\arraybackslash}p{34mm}
 >{\raggedright\arraybackslash}p{34mm}
 >{\raggedright\arraybackslash}p{34mm}@{}}
\toprule\noalign{}
Role & Exact source ID & \(B\) & \(\Lambda\) & \(q\) & \(h\) \\
\midrule\noalign{}
\endhead
\bottomrule\noalign{}
\endlastfoot
Physical numerator & d16\_\allowbreak{}39000 & 0.96934050457 & 3.328741896069 & 0.572847762211 & 0.054998999999 \\
Marked HIGH count & d16\_\allowbreak{}42000 & 0.964090187378 & 3.198839752287 & 0.553741151397 & 0.054998999999 \\
\end{longtable}
\endgroup
\end{landscape}

The 64 nonnatural nodes listed below and the proved natural input give these two conclusions.

\hypertarget{sources:premise-induction-rule-and-exact-constants}{%
\subsection{2. Premise, induction rule and exact constants}\label{sources:premise-induction-rule-and-exact-constants}}

The sole initial source premise is

\[
 \theta_I=0.87481891,\qquad
 \kappa_I=2\theta_I-1=0.74963782,\qquad
 K=6\kappa_I=4.49782692,
\]

and the inherited natural theorem \(G(1,K,0.055)\) in the global convention (S1). The analytic induction rule is the global theorem (G1)--(G17), using divisor transfer (1)--(9), corrected amplification (CA1)--(CA14), and HIGH counting (LD1)--(LD17). The finite schedule uses elementary centered saving \(\alpha=499/1000\), from the contour \(501/1000\).

Every decimal in this appendix is an exact rational number. The fixed accounting parameters are

\[
 r=10^{-8},\quad e_H=10^{-7},\quad
 \mu=\tau=p=\eta=10^{-12},\quad r_0=1/50,\quad
 \alpha=499/1000.
\tag{S2}
\]

Here \(\mu\) is the width mesh allowance, \(\tau\) the tiny-slot allowance, \(p\) the radius allowance, and \(\eta\) the analytic allowance. The actual analytic choices are made sufficiently small to stay below these allowances, after the finite positive reserves have been measured.

For a row of the table with data \(B,\Lambda,q,h\) and parent \(B_0,\Lambda_0,h_0\), define the exact crossing

\[
 b_\times=\frac{B_0-B}{\Lambda_0-\Lambda}.
\tag{S3}
\]

The verifier checks \(0\le b_\times\le\min(h,h_0)\) and \(\Lambda<\Lambda_0\). For \(0\le s\le b_\times\), the new face is no larger than the parent face, so the parent theorem proves it. On the remaining interval \([b_\times,h]\), fix \(q\) as in the table, and at each width-cell upper endpoint \(b\) use

\[
\begin{aligned}
t&=B-\Lambda(b-\mu),& T&=t+b,& S&=b-\mu-\tau,\\
a&=\frac{q\{b+(t-B_0)/\Lambda_0\}+r}
        {1-6qB_0/\Lambda_0},&
F&=1-6a,& F_p&=F-p,\\
L_{\rm ref}&=\frac{49}{100}-Kb,&
z_0&=\frac{B_0F-t}{\Lambda_0}.
\end{aligned}
\tag{S4}
\]

Thus \(a-q\{b+(t-B_0F)/\Lambda_0\}=r\) exactly. All scalar quantities are affine in \(b\).

The scalar certificate checks the positivity of

\[
\begin{gathered}
1-6qB_0/\Lambda_0,\ S,\ a,\ 1/75-a,\ F-T/2-4p,\\
B_0F-t-b-4p,\quad 1-qh_0,\\
r-(qB_0/\Lambda_0+qh_0+4)p-10\eta,\quad
a-10p-10\eta,\quad L_{\rm ref},\\
1-L_{\rm ref}-1/2-Kb-10p,\\
5/6+\alpha L_{\rm ref}-T-9a-e_H,\\
1-T-5a-e_H,\quad 1-r_0-T-4a-e_H,\\
1-r_0-T-(9/2)a-e_H,\quad L_{\rm ref}-7a-e_H,
\end{gathered}
\tag{S5}
\]

and the weak guard \(h_0F-z_0\ge0\). Equality in the latter is permitted because the capped LOW branch retains the positive reserve in (S5). Both width endpoints are checked.

For \(q\le\nu\le3/4\), put

\[
\begin{aligned}
\rho&=1-\frac{10\nu}{3(5-2\nu)},&
\phi&=\frac{\nu(5-6\nu)}{5-2\nu},\\
A_\nu&=\frac{14\nu(5-6\nu)}{115-54\nu},&
D_\nu&=\frac{2\nu(5-6\nu)^2}
 {3(5-2\nu)(115-54\nu)},\\
C_\nu&=\frac{16\nu(5-6\nu)^2}
 {3(5-2\nu)(185-138\nu)},&
H&=1-\rho-\nu T-a(5-6\rho).
\end{aligned}
\tag{S6}
\]

The six density gates are

\[
\begin{gathered}
H+\phi S>e_H,\qquad
H+A_\nu S+D_\nu F_p>e_H,\qquad
H+C_\nu F_p>e_H,\\
H+\nu F_p/2-(1-\nu)b>e_H,\\
H+(\rho-1/6)F_p-(1-\nu)b>e_H,\\
H+\nu(F_p-r_0)/2-(1-\rho)r_0-(1-\nu)b>e_H.
\end{gathered}
\tag{S7}
\]

The verifier also proves the common-width rate \(\nu-\phi-2C_\nu>0\) on the entire interval \([1/2,3/4]\). All scheduled \(q\)'s lie in that interval. Multiplication by the positive denominators gives polynomials of degree at most three in \(\nu\). Exact Bernstein positivity at both \(b\) endpoints proves the full continuous width-density region. This is a continuum certificate; no density mesh is used in the proof.

\hypertarget{sources:complete-acyclic-parameter-table}{%
\subsection{3. Complete acyclic parameter table}\label{sources:complete-acyclic-parameter-table}}

The table is in dependency order. Each nonnatural parent appears strictly earlier. The natural row is the only premise. The six columns below are the complete numerical inputs needed to reproduce the source certificate.

\begin{landscape}
\begingroup % exact-table-layout
\fontsize{7.5}{9.5}\selectfont\setlength{\tabcolsep}{2pt}\renewcommand{\arraystretch}{1.18}
\begin{longtable}{@{}
 >{\raggedright\arraybackslash}p{42mm}
 >{\raggedright\arraybackslash}p{42mm}
 >{\raggedright\arraybackslash}p{34mm}
 >{\raggedright\arraybackslash}p{34mm}
 >{\raggedright\arraybackslash}p{34mm}
 >{\raggedright\arraybackslash}p{34mm}@{}}
\toprule\noalign{}
Source ID & Parent ID & \(B\) & \(\Lambda\) & \(q\) & \(h\) \\
\midrule\noalign{}
\endhead
\bottomrule\noalign{}
\endlastfoot
natural & Natural premise & 1 & 4.49782692 & --- & 0.055 \\
g1\_\allowbreak{}14000 & natural & 0.999728337095 & 4.477808881865 & 0.746594775892 & 0.029031830945 \\
g2\_\allowbreak{}14000 & g1\_\allowbreak{}14000 & 0.999594902255 & 4.468137794316 & 0.745036095459 & 0.02928436799 \\
g3\_\allowbreak{}16000 & g2\_\allowbreak{}14000 & 0.998823044164 & 4.41636472517 & 0.737187559984 & 0.030553384982 \\
g4\_\allowbreak{}16000 & g3\_\allowbreak{}16000 & 0.998440012556 & 4.391616116947 & 0.733210163288 & 0.031189418935 \\
g5\_\allowbreak{}18000 & g4\_\allowbreak{}16000 & 0.997436711696 & 4.331747952564 & 0.724157400173 & 0.032632588034 \\
g6\_\allowbreak{}18000 & g5\_\allowbreak{}18000 & 0.996942367778 & 4.303350744585 & 0.719597099949 & 0.033351088663 \\
g7\_\allowbreak{}20000 & g6\_\allowbreak{}18000 & 0.995776809752 & 4.241113549589 & 0.710258669042 & 0.034819680228 \\
g8\_\allowbreak{}20000 & g7\_\allowbreak{}20000 & 0.995208400142 & 4.211810110132 & 0.70555065255 & 0.035550781757 \\
g9\_\allowbreak{}20000 & g8\_\allowbreak{}20000 & 0.994935006182 & 4.197944457713 & 0.703317368738 & 0.035896662305 \\
g9\_\allowbreak{}22000 & g8\_\allowbreak{}20000 & 0.993898618067 & 4.14861000466 & 0.696148985463 & 0.037010817529 \\
g10\_\allowbreak{}20000 & g9\_\allowbreak{}20000 & 0.994804310728 & 4.191367962573 & 0.702256856613 & 0.036060700065 \\
g10\_\allowbreak{}22000 & g9\_\allowbreak{}22000 & 0.993266628247 & 4.119077756608 & 0.691397158393 & 0.037738686793 \\
g11\_\allowbreak{}20000 & g10\_\allowbreak{}20000 & 0.994741992764 & 4.188245023928 & 0.701752989642 & 0.036138585534 \\
g11\_\allowbreak{}24000 & g10\_\allowbreak{}22000 & 0.991817177568 & 4.055302505813 & 0.681991012399 & 0.039182743845 \\
g12\_\allowbreak{}22000 & g11\_\allowbreak{}20000 & 0.993683145537 & 4.137770984837 & 0.694394223566 & 0.037282481135 \\
g12\_\allowbreak{}24000 & g11\_\allowbreak{}24000 & 0.99112521825 & 4.025739050372 & 0.677222297508 & 0.039903896403 \\
g13\_\allowbreak{}22000 & g12\_\allowbreak{}22000 & 0.993174584155 & 4.114139074471 & 0.69058794895 & 0.037864981859 \\
g13\_\allowbreak{}24000 & g12\_\allowbreak{}24000 & 0.990798325415 & 4.011960981086 & 0.674992876714 & 0.040240282621 \\
g14\_\allowbreak{}22000 & g13\_\allowbreak{}22000 & 0.992932619959 & 4.103028665954 & 0.688794396301 & 0.038138921239 \\
g14\_\allowbreak{}26000 & g13\_\allowbreak{}24000 & 0.989389157004 & 3.955320276699 & 0.666819652435 & 0.041483679205 \\
g15\_\allowbreak{}24000 & g14\_\allowbreak{}22000 & 0.991675436417 & 4.048135087056 & 0.680815207376 & 0.039365150407 \\
g15\_\allowbreak{}26000 & g14\_\allowbreak{}26000 & 0.988724970081 & 3.929254639921 & 0.662597284058 & 0.042114349434 \\
g16\_\allowbreak{}24000 & g15\_\allowbreak{}24000 & 0.991076914713 & 4.022653080472 & 0.676701859858 & 0.039986846052 \\
g16\_\allowbreak{}28000 & g15\_\allowbreak{}26000 & 0.987041667591 & 3.866371669216 & 0.653500312943 & 0.043484255538 \\
g17\_\allowbreak{}24000 & g16\_\allowbreak{}24000 & 0.990794489129 & 4.010768775632 & 0.674778285668 & 0.040277012093 \\
g17\_\allowbreak{}28000 & g16\_\allowbreak{}28000 & 0.9862556267 & 3.837720344138 & 0.6488403297 & 0.044172852358 \\
g18\_\allowbreak{}26000 & g17\_\allowbreak{}24000 & 0.989406842387 & 3.955032836539 & 0.666751225982 & 0.041498424982 \\
g18\_\allowbreak{}30000 & g17\_\allowbreak{}28000 & 0.984404340698 & 3.773326953405 & 0.639576460655 & 0.045556254073 \\
g19\_\allowbreak{}26000 & g18\_\allowbreak{}26000 & 0.9887529443 & 3.929379290092 & 0.662595564572 & 0.042119111023 \\
g19\_\allowbreak{}30000 & g18\_\allowbreak{}30000 & 0.983549291104 & 3.744274316547 & 0.634824648122 & 0.046251281505 \\
g20\_\allowbreak{}26000 & g19\_\allowbreak{}26000 & 0.988447215192 & 3.917514278132 & 0.660667730523 & 0.042406535592 \\
g20\_\allowbreak{}30000 & g19\_\allowbreak{}30000 & 0.983157119711 & 3.731089626457 & 0.632659212628 & 0.046567475664 \\
d1\_\allowbreak{}27000 & g20\_\allowbreak{}26000 & 0.987640126735 & 3.886935287304 & 0.656240492195 & 0.043073258508 \\
d1\_\allowbreak{}31000 & g20\_\allowbreak{}30000 & 0.982201211907 & 3.699630048938 & 0.628164141304 & 0.047234056529 \\
d2\_\allowbreak{}28000 & d1\_\allowbreak{}27000 & 0.986568011422 & 3.847553285526 & 0.650405110191 & 0.043947667353 \\
d2\_\allowbreak{}32000 & d1\_\allowbreak{}31000 & 0.980956443971 & 3.659765312852 & 0.622307764492 & 0.048097948757 \\
d3\_\allowbreak{}29000 & d2\_\allowbreak{}28000 & 0.985344081992 & 3.8040695482 & 0.643920667613 & 0.044914128421 \\
d3\_\allowbreak{}32500 & d2\_\allowbreak{}32000 & 0.979976585481 & 3.629051065709 & 0.617602569731 & 0.04878594477 \\
d4\_\allowbreak{}30000 & d3\_\allowbreak{}29000 & 0.984023104139 & 3.758696083224 & 0.637142847271 & 0.045919661693 \\
d4\_\allowbreak{}33000 & d3\_\allowbreak{}32500 & 0.979113180459 & 3.602471114575 & 0.613577421569 & 0.049374470046 \\
d5\_\allowbreak{}31000 & d4\_\allowbreak{}30000 & 0.982631740443 & 3.712471468246 & 0.630238005469 & 0.046939764667 \\
d5\_\allowbreak{}34000 & d4\_\allowbreak{}33000 & 0.977860651719 & 3.564825264904 & 0.608110688495 & 0.050176812956 \\
d6\_\allowbreak{}32000 & d5\_\allowbreak{}31000 & 0.981183111871 & 3.665885657775 & 0.623284086622 & 0.047963168636 \\
d6\_\allowbreak{}34500 & d5\_\allowbreak{}34000 & 0.976856579435 & 3.535230090578 & 0.603586734106 & 0.050834011243 \\
d7\_\allowbreak{}32500 & d6\_\allowbreak{}32000 & 0.980108998561 & 3.63215533704 & 0.618084793667 & 0.048722743123 \\
d7\_\allowbreak{}35500 & d6\_\allowbreak{}34500 & 0.975502114374 & 3.496250704498 & 0.597915939364 & 0.051662290537 \\
d8\_\allowbreak{}33000 & d7\_\allowbreak{}32500 & 0.979202261989 & 3.604218839864 & 0.613836558219 & 0.049343482376 \\
d8\_\allowbreak{}36000 & d7\_\allowbreak{}35500 & 0.974436908062 & 3.46617658788 & 0.593307531891 & 0.052328493983 \\
d9\_\allowbreak{}34000 & d8\_\allowbreak{}33000 & 0.977929649618 & 3.56595586436 & 0.608268400998 & 0.050160409949 \\
d9\_\allowbreak{}36500 & d8\_\allowbreak{}36000 & 0.973496789612 & 3.440059281266 & 0.589358458856 & 0.052899909767 \\
d10\_\allowbreak{}34500 & d9\_\allowbreak{}34000 & 0.976916338897 & 3.536083748735 & 0.60369906868 & 0.050824121205 \\
d10\_\allowbreak{}37500 & d9\_\allowbreak{}36500 & 0.972122320321 & 3.402690702652 & 0.583975777317 & 0.053683456432 \\
d11\_\allowbreak{}35500 & d10\_\allowbreak{}34500 & 0.975557513138 & 3.49697622457 & 0.598007471598 & 0.051655367964 \\
d11\_\allowbreak{}38000 & d10\_\allowbreak{}37500 & 0.971029173485 & 3.373491761056 & 0.579503754483 & 0.054326688134 \\
d12\_\allowbreak{}36000 & d11\_\allowbreak{}35500 & 0.974490243582 & 3.466842950752 & 0.593389683168 & 0.052322908519 \\
d12\_\allowbreak{}38500 & d11\_\allowbreak{}38000 & 0.970056014935 & 3.347887196314 & 0.575636335117 & 0.054883739096 \\
d13\_\allowbreak{}37000 & d12\_\allowbreak{}36000 & 0.973070357376 & 3.427670025232 & 0.587700845567 & 0.053150754524 \\
d13\_\allowbreak{}39500 & d12\_\allowbreak{}38500 & 0.968618636223 & 3.310830868376 & 0.570323433367 & 0.054998999999 \\
d14\_\allowbreak{}37500 & d13\_\allowbreak{}37000 & 0.971961487485 & 3.397636928241 & 0.583092034117 & 0.053814209892 \\
d14\_\allowbreak{}40500 & d13\_\allowbreak{}39500 & 0.966952592645 & 3.26886140552 & 0.564180496021 & 0.054998999999 \\
d15\_\allowbreak{}38500 & d14\_\allowbreak{}37500 & 0.970487094978 & 3.358580552439 & 0.577434117401 & 0.054635093269 \\
d15\_\allowbreak{}41000 & d14\_\allowbreak{}40500 & 0.965705536457 & 3.237999482351 & 0.559400373933 & 0.054998999999 \\
d16\_\allowbreak{}39000 & d15\_\allowbreak{}38500 & 0.96934050457 & 3.328741896069 & 0.572847762211 & 0.054998999999 \\
d16\_\allowbreak{}42000 & d15\_\allowbreak{}41000 & 0.964090187378 & 3.198839752287 & 0.553741151397 & 0.054998999999 \\
\end{longtable}
\endgroup
\end{landscape}
